\documentclass[11pt,letterpaper]{article}
\usepackage[T1]{fontenc}
\usepackage[utf8]{inputenc}
\usepackage{lmodern}
\pdfmapfile{+lm.map}
\pdfmapfile{+cm.map}
\pdfmapfile{+symbols.map}
\usepackage[margin=1in]{geometry}
\usepackage{amsmath,amssymb,amsthm,mathtools}
\usepackage{microtype,booktabs,tabularx,array,xcolor,listings,fancyhdr}
\usepackage[hidelinks]{hyperref}
\definecolor{ink}{HTML}{18344D}
\hypersetup{pdftitle={iPRA with evaluation at witness discharge},
 pdfauthor={Research prepared for Vladimir Reshetnikov},
 pdfsubject={Three inference families on raw terms, exact bi-interpretability, and complete checkers}}
\newtheorem{theorem}{Theorem}[section]
\newtheorem{lemma}[theorem]{Lemma}
\newtheorem{proposition}[theorem]{Proposition}
\newtheorem{corollary}[theorem]{Corollary}
\theoremstyle{definition}\newtheorem{definition}[theorem]{Definition}
\theoremstyle{remark}\newtheorem{remark}[theorem]{Remark}
\setlength{\parskip}{3pt}
\setlength{\emergencystretch}{3em}
\setlength{\headheight}{14pt}
\pagestyle{fancy}\fancyhf{}
\fancyhead[L]{\small iPRA: evaluation at witness discharge}
\fancyhead[R]{\small Research note}
\fancyfoot[C]{\thepage}
\lstset{language=Python,basicstyle=\ttfamily\scriptsize,
 columns=fullflexible,keepspaces=true,breaklines=true,
 showstringspaces=false,frame=single,rulecolor=\color{black!20},
 xleftmargin=4pt,xrightmargin=4pt,aboveskip=9pt,belowskip=9pt}
\setcounter{tocdepth}{1}
\title{\color{ink}iPRA with evaluation at witness discharge\\[7pt]
\large Three inference families on raw terms}
\author{Research prepared for Vladimir Reshetnikov}
\date{November 11, 2026}
\begin{document}
\maketitle
\begin{abstract}
We give a three-family presentation of full first-order intuitionistic
primitive recursive arithmetic on raw scalar-iterator terms. The primitive
families are clause use, paired introduction, and witness discharge with
one elementary evaluation; universal reflexivity is the one explicit fixed
axiom clause. There is no numerical normalization. The evaluation map
contains the two existing iterator-evaluation patterns and identity as its
fallback. We prove exact equivalence, on every open end sequent, with the
reference calculus having separate evaluation and witness elimination.
The inherited whole-domain interpretations with iPRA therefore retain
actual equality and identity numerical composites. Both complete formula
round trips and equiconsistency are included.

The complete raw checker has 1,892 normalized bytes, 643 significant tokens
and 695 AST nodes, versus 1,995, 665 and 717 for its matched control.
A zero-normal companion has 1,858 bytes, 633 tokens and 697 AST nodes.
Both certificate predicates are primitive recursive without fuel.
The improvement is a local rule fusion and a measured reduction in the
stated checker family, not an absolute minimum or a claim to surpass the
repository's different compact codecs. The archive includes complete
checkers, proof translations, independent replay and the mathematical
foundations needed for the bi-interpretation.
\end{abstract}
\clearpage
\tableofcontents
\clearpage
\section{The result and its comparison target}
\label{sec:result}

The proposed simplification is a calculus with \textbf{three inference
families on completely raw terms}: use a retained clause, introduce a
paired formula, and discharge a witness while performing one elementary
evaluation. There is one explicit initial axiom clause, universal
reflexivity. Numerical computation is never an implicit conversion rule.
Neither full iterator normalization nor zero normalization is needed in
this primary presentation.

The construction is compared with ProveIt's \texttt{ConciseIPRA} at commit
\texttt{11814771b749}. The inspected frontier
includes the seventy-three filed packages and all six archives then present
in \texttt{docs/incoming}. In particular, the comparison includes
\texttt{witness\_discharge}, \texttt{clause\_discharge},
\texttt{deterministic\_opening}, \texttt{first\_instance},
\texttt{inferred\_witness}, \texttt{prefix\_clauses}, and
\texttt{substitutive\_discharge}. The last group matters: a proposal that
only removes the supplied equality proof or replacement term from the
earlier Name rule would already be present there.

What is new here is the fusion of the existing \emph{computation-and-equality
evaluation rule} with witness elimination. In contrast, the earlier
three-family Name presentations put both iterator equations into numerical
normal forms. The incoming raw contextual-discharge variant retains a
separate computation family. The present raw calculus has three families
while keeping ordinary, nonnormalizing substitution.

An optional companion uses zero-normal terms and the existing single
SeedEval pattern. It also has three families, and has the smaller complete
checker. The raw variant is the cleaner recommendation when the priority
is the trusted mathematical specification; the companion is preferable
when normalized checker size is decisive. Both preserve the entire
intuitionistic first-order theory, not merely its equations.

The result is a local simplification of the current presentation. It does
not prove an absolute minimum over all axiomatizations, encodings or
programming languages. Nor does it beat the repository's separate
508-byte compact codec, whose proof representation and validation boundary
are different. The whole-program comparison in Section~\ref{sec:metrics}
makes these boundaries explicit.

\section{The complete raw presentation}
\label{sec:raw}

\subsection{Syntax, clauses, and the one fixed axiom}

There is one numerical sort. Raw terms and native formulas are
\[
 t::=x\mid 0\mid S(t)\mid I(n,a;u.h),
 \qquad A::=(s=t)\mid D_{x,y}A.
\]
Only $u$ in $h$ is bound by the iterator. The two variables of $D$ are
bound in its matrix. Bound variables are renamed before capture-avoiding
substitution; the executable representation instead uses bound indices.
All finite expressions of this grammar are terms or formulas: there is
no normal-form side condition and no identification by computation.
In particular, $I(0,a;u.h)$ and $a$ are distinct terms.

A reusable clause has the form $[\vec A;y.B(y)]$, where $\vec A$ is empty
or consists of one formula, and $y$ is absent from that antecedent.
Its ordinary reading is, respectively, $\forall y\,B(y)$ or
$A\to\forall y\,B(y)$. A context is a finite list of clauses. The
abbreviation $\Gamma,A$ means addition of the unconditional constant
clause $[\,;v.A]$, with $v$ unused.

Every arithmetic proof begins with the explicitly counted clause
\[
                  \boxed{\mathcal R=[\,;v.v=v].}
\]
It supplies equality reflexivity by clause use. The native contradiction
is the formula $0=S0$. There are no other primitive axiom clauses.

\subsection{One deterministic elementary evaluation}

For arbitrary formulas $E,F$, define $\rho_E(F)$ by the following three
disjoint cases. Every match is literal syntax, including equality of the
displayed iterator bodies after the chosen representation of binders.
\begin{align}
 \rho_{s=t}\bigl(I(0,s;u.h)=r\bigr)&=(r=t),
       \tag{Z}\label{eq:rho-zero}\\
 \rho_{I(c,a;u.h)=t}\bigl(I(Sc,a;u.h)=r\bigr)
       &=\bigl(r=h[t/u]\bigr),
       \tag{S}\label{eq:rho-succ}\\
 \rho_E(F)&=F\qquad\text{in every other case}.
       \tag{Id}\label{eq:rho-id}
\end{align}
The zero case has a literal zero counter; the successor case has a
successor constructor, so they cannot overlap. In the second case the
entire predecessor iterator on the left of $E$ must match the one obtained
from the left of $F$ by deleting its counter's outer successor.
The replacement term, body, and counter are inferred from the two formulas;
no term template or computation description is separate rule data.

The substitution $h[t/u]$ is a single ordinary syntactic substitution.
The result is not normalized. An auxiliary formula $E$ with a $D$ root
always selects Id. The identity case does not derive $F$ without a premise:
the rule below explicitly requires a proof of $F$.

\subsection{All three inference families}

\paragraph{Use.}
For any retained clause $[\vec A;y.B]\in\Gamma$ and term $t$,
\[
 \frac{\{\Gamma\vdash A:A\in\vec A\}}
      {\Gamma\vdash B[t/y]}.
 \tag{Use}\label{eq:raw-use}
\]
There is no proof premise for an unconditional clause, and one for a
conditional clause. This is the only leaf mechanism; reflexivity is an
instance of Use at $\mathcal R$.

\paragraph{Introduce.}
\[
 \frac{\Gamma,B(t,\eta)\vdash B(t,Sz)}
      {\Gamma\vdash D_{x,y}B(x,y)},\qquad
 \eta=\begin{cases}
 z,&B\text{ is literally an equation and }t\equiv0,\\
 0,&\text{otherwise}.
 \end{cases}
 \tag{Intro}\label{eq:raw-intro}
\]
Here $z$ is fresh for $\Gamma$, the witness $t$, and the passive parameters
of $B$. The witness test is on raw syntax. The progress case is restricted
to literal equations; it does not introduce induction on quantified formulas.

\paragraph{Evaluate and discharge.}
For a fresh numerical parameter $b$, write
$\mathcal B_b=[B(b,0);y.B(b,y)]$. The third family is
\[
 \frac{\Gamma\vdash D_{x,y}B\qquad
       \Gamma,\mathcal B_b\vdash E\qquad
       \Gamma,\mathcal B_b\vdash F}
      {\Gamma\vdash\rho_E(F)}.
 \tag{Discharge}\label{eq:raw-discharge}
\]
The parameter $b$ is fresh for $\Gamma$ and the passive parameters of the
major formula. Both branch proofs use the same extended context. The
computed conclusion must be well scoped in the outer numerical scope;
equivalently, it has no free occurrence of $b$. The intermediate formulas
$E,F$ may contain $b$.

These are the entire primitive rules of the calculus, denoted $\mathcal E_3$.
There is no separate hypothesis, equality transport, computation,
elimination, case, induction, normalization, weakening or cut instruction.
The two nonidentity patterns inside $\rho$ remain explicitly counted in
the specification. Calling this ``three families'' does not conceal their
existence or establish a lower bound.

\section{Exact equivalence on raw open sequents}
\label{sec:raw-equivalence}

Let $\mathcal V_4$ be the same syntax, clauses, Use and Intro, with ordinary
paired elimination and the following two-case Eval family in place of
Discharge:
\[
 \frac{\Gamma\vdash s=t\quad\Gamma\vdash I(0,s;u.h)=r}
      {\Gamma\vdash r=t},\qquad
 \frac{\Gamma\vdash I(c,a;u.h)=t\quad
       \Gamma\vdash I(Sc,a;u.h)=r}
      {\Gamma\vdash r=h[t/u]}.
\]
Its ordinary paired elimination is
\[
 \frac{\Gamma\vdash D_{x,y}B\quad
       \Gamma,\mathcal B_b\vdash A}
      {\Gamma\vdash A},\qquad b\notin\operatorname{FV}(A).
\]
This is the repository's raw uniform-clause calculus: Eval is the inherited
two-case arithmetic interface, and $\mathcal R$ replaces a separate
reflexivity rule. The pinned Lean files \texttt{NativeEval.lean},
\texttt{NativeClausal.lean}, and \texttt{NativeUniformClauses.lean}
establish the corresponding reference equivalences~\cite{repo,rawformal}.
The new fusion theorem below is proved here, rather than attributed to
those existing formalizations.

\begin{lemma}[Structural operations]
\label{lem:raw-structural}
Weakening and injective renaming of free parameters are admissible in
$\mathcal E_3$ and $\mathcal V_4$.
\end{lemma}
\begin{proof}
Induct on derivations, freshening a bound eigenparameter before adding an
outer parameter or clause. Injective renaming preserves and reflects the
literal matches in $\rho$, preserves the zero/successor constructors,
commutes with capture-avoiding substitution, and preserves the Intro mode.
Weakening changes none of that data. No assertion that arbitrary numerical
substitution preserves an instruction's selected case is needed.
\end{proof}

\begin{lemma}[A compound theorem with no arithmetic premise]
\label{lem:theta}
For every clause context, $\Theta=D_{x,y}(0=0)$ is derivable in both systems,
even if that context does not contain $\mathcal R$.
\end{lemma}
\begin{proof}
Introduce at witness $0$. The temporary progress assumption and the
successor target are both $0=0$. Use that newly added unconditional
assumption. Thus the proof uses only Intro and Use.
\end{proof}

\begin{theorem}[Identity on the complete consequence relation]
\label{thm:raw-equivalence}
For every well-scoped raw clause context $\Gamma$ and formula $A$,
\[
               \Gamma\vdash_{\mathcal E_3}A
               \quad\Longleftrightarrow\quad
               \Gamma\vdash_{\mathcal V_4}A.
\]
Both translations are effective transformations of finite proof trees and
preserve the entire end sequent literally. The statement does not require
$\mathcal R$ to be present.
\end{theorem}
\begin{proof}
From left to right, translate a Discharge node's three premises. If its
$\rho$ computation selects Z or S, apply the corresponding Eval inside the
witness scope. If it selects Id, the translated branch proof of $F$ already
has the required conclusion. Apply ordinary paired elimination. Its
eigenparameter condition is exactly the outer-scope condition of Discharge.
Use and Intro translate without change.

Conversely, translate a reference elimination node by Discharge. Use as
auxiliary $E$ the same compound major $D_{x,y}B$, weakened into the fresh
branch. Since $E$ is compound, $\rho_E(A)=A$, so the original conclusion
is unchanged.

Translate either Eval node by Discharge with the fixed major proof of
$\Theta$ from Lemma~\ref{lem:theta}. Weaken its original equation premises
into that fresh scope, using them as $E,F$. The intended Z or S pattern
matches and reproduces the Eval conclusion. All data of that conclusion
are from the outer scope. Recursing over the finite tree, with fresh
parameters at every inserted scope, proves the two transformations.
\end{proof}

\paragraph{Direct ordinary soundness.}
Read $D_{x,y}B$ as
$\exists x(B(x,0)\to\forall y\,B(x,y))$ and read equality as actual
equality. Eliminate the existential at $b$. In case Z, the zero computation
equation and $s=t$ turn the second branch equation into $r=t$.
In case S, successor computation and the predecessor equation give
$r=h[t/u]$. Id simply retains the branch formula. Each conclusion can
leave the witness scope because its free occurrence of $b$ is forbidden.
This is intuitionistic reasoning and does not use a truth-table argument
or a restriction to the standard natural numbers.

\section{Where equality, computation and induction come from}
\label{sec:bootstrap}

The exact comparison theorem already transfers completeness. The following
derivations expose the arithmetic content and show that no equality law is
being silently supplied by the checker. Write
$J_s=I(0,s;u.0)$. All Eval applications below can be expanded by
Theorem~\ref{thm:raw-equivalence} into the three-family calculus.

\paragraph{Joining equations.}
From $s=t$ and reflexivity of $J_s$, Z yields $J_s=t$.
From $s=r$ and that equation, a second Z yields $t=r$.
Consequently $s=t,s=r\vdash t=r$. Reflexivity supplies symmetry, and
joining with symmetry supplies transitivity.

\paragraph{Both computation equations.}
Z on reflexivity of $a$ and of $I(0,a;u.h)$ yields
$I(0,a;u.h)=a$. S on reflexivity of the predecessor and successor
iterators yields
\[
           I(Sc,a;u.h)=h[I(c,a;u.h)/u].
\]
These are theorems. They are not reductions performed by raw substitution.

\paragraph{Transport in any term.}
Suppose $s=t$ and $h(s)=r$. Put
$J=I(0,s;u.h)$ and $K=I(S0,s;u.h)$.
Zero computation gives $J=s$; successor computation and that equation
give $K=h(s)$, using S. Joining gives $K=r$.
The zero Eval on $s=t$ and reflexivity of $J$ gives $J=t$.
The successor Eval on $J=t$ and $K=r$ now yields $r=h(t)$.
This is the inherited crossed transport rule in an arbitrary term template.
It supplies congruence and transport in an equation. Transport through a
compound formula follows by external induction on its $D$ depth, opening
the old $D$, introducing the new $D$ at the fresh nonzero-syntax witness,
and transporting its base and successor instances.

\paragraph{Equation induction.}
For a literal equation $E(y)$, choose an unused witness coordinate $x$.
The progress instance of Intro at witness $0$ turns
$\Gamma,E(z)\vdash E(Sz)$ into $D_{x,y}E(y)$. Derived ordinary elimination
opens the clause $[E(0);y.E(y)]$. A supplied base proof of $E(0)$ and Use
at any term $t$ derive $E(t)$. This is contextual induction on equations.
No induction on compound formulas has been added.

\paragraph{The case mode omitted at a literal zero witness.}
It is recovered at the raw witness $w=I(0,0;u.0)$, which is not literally
$0$ and is provably equal to $0$. Transport converts the local assumption
and successor target between the witnesses $w$ and $0$. Ordinary-assumption
cut, proved by induction on derivations after freshening parameters, inserts
the original case proof. Thus the single introduction recovers both
reference introductions. This argument uses raw syntax, and is one reason
not to silently identify computed terms in the primary calculus.

The full native logical interface and the arithmetic bridge are recorded
in Sections~\ref{sec:interpretations} and the attributed appendices. The
logical macros, predecessor construction, equality reflection and pairing
compiler are inherited; the new theorem changes their proof interface
without changing their conclusions.

\section{The smaller zero-normal companion}
\label{sec:zero}

There is a second presentation, $\mathcal E^0_3$, with the same three
inference families and the same fixed axiom clause. Its terms are canonical
only for the deletion
\[
                         I(0,a;u.h)\longrightarrow a.
\]
Let $N_0$ perform this deletion hereditarily. Every rule substitution is
followed by $N_0$. A successor-headed iterator is not unfolded by this
syntactic operation. The evaluation map is instead
\[
 \sigma_E(F)=
 \begin{cases}
  r=N_0\bigl(h[I(c,t;u.h)/u]\bigr),
    &E\equiv(s=t),\quad F\equiv(I(Sc,s;u.h)=r),\\[3pt]
  F,&\text{otherwise}.
 \end{cases}
\]
The initializer in the equation $E$ must match the initializer of the
iterator in $F$. The previous iterate is constructed from the replacement
initializer $t$; it is not an additional premise. This is precisely the
existing SeedEval rule, executed only at witness discharge.

Use and Intro are those of Section~\ref{sec:raw}, with normalized
substitution and the literal witness test on canonical syntax. Discharge
exports $\sigma_E(F)$ instead of $\rho_E(F)$, subject to the same outer
scope condition. There is one arithmetic pattern and the identity case.

The proof of Theorem~\ref{thm:raw-equivalence} applies verbatim with
SeedEval in place of the two raw Eval cases. Its separate-rule control has
Use, Intro, Open and SeedEval. The compound major still recovers Open, and
$\Theta$ still recovers SeedEval. In the reverse simulation no arbitrary
substitution into an existing derivation is assumed.

\paragraph{The arithmetic interface is unchanged.}
At counter $0$, SeedEval reads
\[
             a=b,\ I(S0,a;u.h)=r\quad\Longrightarrow\quad r=N_0(h[b/u]).
\]
Taking $h=u$ and using two such evaluations derives the joining rule for
equations. Taking arbitrary $h$ derives transport by the same two-step
pivot construction. Two reflexivities give successor computation.
Zero computation is canonical syntax. This interface is formalized in the
pinned \texttt{NativeZeroSeed.lean} as \texttt{NVPrf}; the equivalence
\texttt{nvprf\_iff\_pra} is already in that source~\cite{zeroformal}.
The clause representation and its fixed reflexivity clause are formalized
in \texttt{NativeZeroClausal.lean}. The present fusion is a new written
equivalence on top of those interfaces, not a claim that its new checker
has been verified by Lean.

\paragraph{Only structural normalization is used.}
Recursively normalize an iterator's children. If its normalized counter is
$0$, return its normalized initializer; otherwise rebuild the iterator.
Every contraction deletes nodes, and none duplicates a body. Structural
induction proves termination, idempotence and substitution coherence
\[
              N_0(N_0(t)[\tau])=N_0(t[\tau]).
\]
Thus this companion does not reinstate the full computation normalizer
removed from the main presentation.

\section{Full first-order iPRA and both interpretation composites}
\label{sec:interpretations}

The reference theory is full first-order intuitionistic primitive recursive
arithmetic: intuitionistic equality logic; a symbol for each fixed finite
primitive-recursive description and its defining equations; separation of
zero from successors; successor injectivity; and induction for every
quantifier-free formula, with parameters and retained assumptions.
In particular, the target is not merely equational PRA or a negative
fragment of iPRA.

\subsection{The intermediate ordinary scalar theory}

Let $T_I$ be ordinary intuitionistic first-order equality logic over the
staged scalar-iterator signature. Its nonlogical axioms and rules are exactly
the two iterator equations, numerical coverage, separation, and contextual
equation induction:
\begin{gather*}
 I(0,a;u.h)=a,\qquad I(Sn,a;u.h)=h[I(n,a;u.h)/u],\\
 \forall x\,(x=0\vee\exists y\,x=Sy),\qquad (0=S0)\to\bot,\\
 \frac{\Gamma\vdash E(0)\qquad\Gamma,E(z)\vdash E(Sz)}
      {\Gamma\vdash E(t)}\quad(E\text{ an equation},\ z\text{ fresh}).
\end{gather*}
The displayed separation sentence is essential to the ordinary realization:
positive equations by themselves do not prohibit a one-element structure.
It is not an extra hidden axiom of the native calculus, where $0=S0$ is
the designated contradiction and explosion is derived.

The native-to-ordinary translation is
\[
 (s=t)^\circ=(s=t),\qquad
 (D_{x,y}B)^\circ=\exists x\bigl(B^\circ(x,0)\to\forall y\,B^\circ(x,y)\bigr).
\]
Unconditional clauses translate to universals and conditional clauses to
implications into universals. All numerical domains are the whole domain,
defined by $x=x$, and equality is actual equality.

\subsection{Explicit logical macros}

Put $C(v,a,b)=I(v,a;u.b)$, with $u$ unused, and
$p(z)=I(z,0;u.C(u,1,0))$. The derived computation equations give the
selector endpoints and $p(1)=1$, $p(2)=0$.
For a formula, define its gate recursively by
\[
 G_v(s=t)=(C(v,0,s)=C(v,0,t)),\qquad
 G_v(D_{x,y}B)=D_{x,y}G_vB,
\]
with fresh bound variables. It is provable at $v=0$ and equivalent to the
original formula at every successor. Formulas can be padded by
$A\dashv\vdash D_{u,y}G_yA$.

The formula selector $M_v(A,B)$ uses $C$ on corresponding sides of two
equations and distributes through aligned $D$ binders. If one formula
is an equation and the other has a $D$ binder, pad the equation using
the preceding gate before continuing. This finite recursive definition,
spelled out in Appendix~\ref{sec:choice-introduction}, gives
$M_0(A,B)\dashv\vdash A$ and $M_{Sv}(A,B)\dashv\vdash B$.
Write $Q_yA=D_{u,y}A$ with $u$ unused. One complete macro inventory is
\[
\begin{aligned}
 \bot_*&=(0=S0),\\
 A\to_*B&=D_{u,y}M_y(A,B),\\
 A\wedge_*B&=D_{u,y}G_y M_{p(y)}(A,B),\\
 \exists_*x\,A(x)&=D_{x,y}G_yA(x),\\
 A\vee_*B&=\exists_*x\,M_x(A,B),\\
 \forall_*x\,A(x)&=A(0)\wedge_* Q_y A(y).
\end{aligned}
\]
Every newly bound variable is fresh, except the explicitly displayed
argument of the formula being bound. These macros have all the usual
intuitionistic introduction and elimination rules. For example,
implication uses the selector's zero row as its assumption and its
successor rows as its conclusion; conjunction is read at the two controls
$1,2$; existential quantification uses the true zero gate. Universal
elimination uses $A(0)$ to activate the $Q$ clause. Under $0=S0$,
selector computation makes any two terms equal, and induction on finite
formula syntax yields explosion.

The appended inherited logical development proves these claims. Its
equivalent universal macro uses a gated paired conjunct; the macro displayed
here is the more direct $A(0)\wedge Q_yA(y)$ also used by the current
Lean development. Both have the ordinary universal reading.

\subsection{The numerical translations}

From scalar terms to iPRA, translate each fixed body $h(u,\vec p)$ by
structural recursion on its finite description. Its iterator is the
primitive-recursive operation
\[
 i_h(0,a,\vec p)=a,\qquad
 i_h(Sn,a,\vec p)=h(i_h(n,a,\vec p),\vec p).
\]
No number is asked to encode a function on which a universal total
evaluator operates.

In the other direction, the inherited scalar development first derives
predecessor from scalar iteration and equation induction. It then proves
successor injectivity, constructs an equality indicator that reflects
actual equality internally, and obtains induction on all quantifier-free
formulas. Pairing and its projections are derived next. For a fixed
primitive recursion
\[
 f(0,\vec p)=g(\vec p),\qquad
 f(Sn,\vec p)=h(n,f(n,\vec p),\vec p),
\]
the scalar state stores a pair consisting of the counter and the current
value. If $L,H$ are its two projections, set
\[
\begin{aligned}
 W(v,\vec p)&=\pi(SL(v),h^*(L(v),H(v),\vec p)),\\
 V(n,\vec p)&=I(n,\pi(0,g^*(\vec p));v.W(v,\vec p)),\\
 f^*(n,\vec p)&=H(V(n,\vec p)).
\end{aligned}
\]
Equation induction proves $L(V(n,\vec p))=n$, and the two recursion
equations for $f^*$ follow. The order of this argument matters: the
predecessor and equality-reflection proofs do not presuppose the general
primitive-recursion compiler. The appended foundations supply that order
and the necessary internal proofs~\cite{choice,bridge,erasure,forward,progress}.

\begin{theorem}[Required bi-interpretability and equiconsistency]
\label{thm:biinterp}
The ordinary realizations of $\mathcal E_3$ and $\mathcal E^0_3$ are
bi-interpretable with full first-order iPRA. Both numerical interpretation
domains are whole, equality is actual equality, and both numerical
composite isomorphisms have graph $x=y$. Both formula composites are
provably equivalent to the identity. Each presentation is equiconsistent
with iPRA.
\end{theorem}
\begin{proof}
Theorem~\ref{thm:raw-equivalence} transfers the full native consequence
relation to the raw uniform-clause reference, whose scalar and iPRA
interpretations are the ones just described. The zero companion has the
same exact comparison with its SeedEval reference.

For the numerical composites, induction on a fixed finite primitive-
recursive description shows that its double translation satisfies the
original defining equations. Internal equation induction proves uniqueness
of that fixed recursion. The analogous induction on scalar syntax proves
the other composite. These are statements for arbitrary internal counters,
not just for standard numeral instances.

The crucial native formula composite is
\[
 D_{x,y}B\dashv\vdash
 \exists_*x\bigl(B(x,0)\to_*\forall_*y\,B(x,y)\bigr).
\]
Forward, open its witness clause, assume its base, and use that clause
at a fresh variable before introducing the ordinary macros. Backward,
eliminate the existential at a fresh witness and use case introduction
with the supplied implication and universal at the successor instance.
The other formula composites follow from the selector and gate identities
and induction on finite formulas. The zero companion adds the internally
provable equivalence of a term with its zero-normal form.

The translations preserve contradiction: ordinary falsity maps to
$0=S0$, and the ordinary reading of $0=S0$ implies falsity by separation.
Effective translations of finite derivations in both directions therefore
give equiconsistency. This proof does not replace the whole domain by a
standard cut, a quotient, or a coded subset.
\end{proof}

The inherited numerical and logical bridge has a Lean formalization in
the pinned repository. The new fusion theorem and the association of its
Python checker with its stated rules are written proofs with executable
validation. No new Lean or Rocq verification is claimed.

\section{Complete certificate checkers}
\label{sec:implementation}

Both delivered acceptance programs are standalone Python modules using only
the standard library. Each reads one JSON value from standard input.
Acceptance means exit status zero and the line \texttt{ACCEPT}. The
document has exactly the keys \texttt{goal} and \texttt{proof}. The
goal is a closed native formula, and the checker infers every intermediate
conclusion from the proof tree. It compares the inferred final formula with
the public goal. There are no trusted proof-producing helpers.

\subsection{The complete public syntax}

The same disjoint-by-sort encoding is used by both programs:
\begin{center}
\begin{tabular}{ll}
\toprule
Object & JSON representation\\
\midrule
Zero & \texttt{0}\\
Free parameter level & a negative integer\\
Bound variable index & a positive integer\\
Successor $St$ & \texttt{[0,t]}\\
Iterator $I(n,a;u.h)$ & \texttt{[n,a,h]}\\
Equation $s=t$ & \texttt{[s,t]}\\
Paired formula $D B$ & \texttt{[B]}\\
\bottomrule
\end{tabular}
\end{center}
An iterator adds one bound index in its last child only; $D$ adds two in
its matrix. At a matrix's outer level, index $1$ is its induction/closure
variable and index $2$ its witness variable. Integers are checked with exact
type tests, so JSON Booleans are not accepted as numerical terms or indices.
Constructor arities and the successor's zero tag are checked as well.

A clause is internally stored as a pair \texttt{(antecedent,body)}.
An empty list in the antecedent position denotes its absence; an actual
formula cannot be an empty list. The initial clause is
\texttt{([], [1,1])}, exactly $\mathcal R$.

The three proof forms are
\begin{lstlisting}[language={}]
["u", clause_index, term, antecedent_proof]
["i", matrix, witness, successor_branch]
["e", major_proof, auxiliary_branch, formula_branch]
\end{lstlisting}
Use ignores its antecedent-proof payload for an unconditional clause. The
payload may be any JSON value in that case; this is an intentional part of
the format, not an unvalidated logical premise. Intro infers the paired
formula from its supplied matrix. Discharge obtains the matrix by checking
its major proof; its two branch conclusions determine the evaluation.
No intermediate branch formula, replacement term or template is separately
supplied to Discharge.

\subsection{Scope and capture invariants}

The proof starts with one clause and no free parameters. If a checking
context contains $n$ clauses, its available free levels are
$1-n,\ldots,-1$. Every Intro or Discharge creates exactly one new numerical
parameter $b=-n$ and adds exactly one clause. This is why clause count can
supply the scope bound in the closed-theorem interface. It is not a claim
that arbitrary mathematical clause contexts determine their free variables.

Every recursive proof result is checked for its formula sort and numerical
scope. A discharge result is checked in the \emph{outer} scope, after its
elementary evaluation. This admits a valid inference in which evaluation
erases the witness from the intermediate formula; it rejects every
inference whose exported formula still contains the witness. Premises
are independently checked in their common local scope before evaluation.

Every term substituted by the implementation is externally scoped and has
no loose positive index. Its internal positive indices are already under
its own binders. Consequently inserting it below an additional binder
requires no index shift. The shared traversal only adjusts the index of
the opened binder and decrements indices above it. In the successor Eval,
the body comes from the root iterator of an equation, so there is precisely
one open iterator variable to instantiate and no hidden surrounding formula
binder. The zero companion has the same invariant.

The checker does not validate an Intro matrix in a separate preliminary
pass. It validates the complete original matrix in the synthesized result
$[B]$ before returning successfully. Thus malformed material cannot become
acceptable merely because branch substitution erases it. The adversarial
suite includes such an attempt in the zero-normal variant.

\subsection{Primitive-recursive checking without fuel}

On standard finite-tree encodings, raw substitution is primitive recursive.
For tree-node size, a conservative bound is
\[
                |\operatorname{sub}(a,t)|
                \le |a|\max(1,|t|).
\]
The traversal visits the original tree and does not traverse newly inserted
copies recursively. A raw successful evaluation performs either no
substitution or one body substitution. The zero companion adds only
structural deletions and at most one constructed predecessor iterator.

Every recursive proof call consumes a proper subcertificate. The finitely
many traversals and substitutions at one node have a fixed polynomial
majorant in the sizes of that node's input data, retained context and
already inferred subproof conclusions. Iterating a generous such majorant
over the finite certificate tree gives primitive-recursive bounds on all
intermediate sizes and verification time. The ideal certificate predicate
is therefore primitive recursive in both presentations, with no additional
fuel field.

This differs from running a full normalizer on arbitrary primitive-
recursive program descriptions. Neither checker asks for a uniform numerical
evaluation of such descriptions or adds a totality assertion for one to
iPRA. Ordinary Python parser, stack and memory limits still bound concrete
runs; they are separate from this ideal finite-tree statement.

\section{Measured comparison and executed evidence}
\label{sec:metrics}

All reported measurements include the entire acceptance module: parsing,
syntax and scope validation, substitution, context operations, inference
checking, final-goal comparison and the acceptance print. Under CPython
3.12.14, remove actual docstrings, apply \texttt{ast.unparse}, and add one
final LF. Count UTF-8 bytes, significant tokens and Python AST nodes.
Gzip uses level $9$ and timestamp zero, and is informational.
The bundled \texttt{measure.py} reproduces the table from the actual source
files. Readable comments and documentation are not part of normalized
source size.

\begin{center}\small
\begin{tabularx}{\textwidth}{@{}Xrrrr@{}}
\toprule
Complete acceptance program & Bytes & Tokens & AST & Gzip9\\
\midrule
Raw evaluation discharge & 1892 & 643 & 695 & 773\\
Matched raw separate rules & 1995 & 665 & 717 & 786\\
Zero-normal evaluation discharge & 1858 & 633 & 697 & 777\\
Matched zero-normal separate rules & 1879 & 645 & 707 & 771\\
Incoming inferred-witness output control & 1983 & 663 & 751 & 812\\
Incoming inferred witness & 1990 & 683 & 771 & 819\\
Incoming explicit clause discharge & 2010 & 636 & 708 & 791\\
\bottomrule
\end{tabularx}
\end{center}

The isolated raw fusion saves 103 bytes, 22 tokens and 22 AST nodes
(1,995 to 1,892 bytes, a 5.16\% reduction). Its gzip size is also 13 bytes
smaller. The isolated zero-normal fusion saves 21 bytes, 12 tokens and
10 AST nodes; its gzip size increases by six bytes. Against the incoming
1,983-byte output-control program, the complete zero companion saves
125 bytes (6.30\%), 30 tokens and 54 AST nodes.


The matched controls use exactly the same syntax and public contract as
their new checkers, with separate Eval/SeedEval and Open instructions.
The difference from those controls isolates the rule fusion. Comparisons
with the incoming full-normalization systems also change the term
discipline and some implementation details; they must not be described as
the isolated saving from fusion.

The zero companion improves all three primary coordinates of the three
listed incoming strict clause-checker programs. The raw checker prioritizes
having no normalization machinery: it has fewer bytes and AST nodes than
those programs, but its token count exceeds the 636-token explicit
clause-discharge program. Neither is asserted to dominate every possible
representation or every checker in the repository. The separate compact
codec in incoming \texttt{deterministic\_opening} is 508 normalized bytes,
and its validated public-goal companion is 780 bytes. They remain smaller
by that coordinate and implement a different certificate interface.

The delivered raw suite translates all 33 pinned upstream raw certificates
into both the fused and separate kernels. It checks 38 accepted examples,
38 exact-goal fused/separate round trips, and 38 reverse expansions in an
independently implemented upstream raw checker. It rejects 34 adversarial
inputs and passes 144 standalone ordinary/optimized-Python CLI checks.

The zero suite replays all 85 pinned SeedEval certificates in their original
checker, translates all 85 into each new kernel, checks both directions
of the generic transformation and 85 fused round trips, and checks six
additional positives and 27 adversaries. Its 66 standalone CLI outcomes
also pass. The package-level reproduction runs both suites from their
packaged paths and separately replays the independent reviewer fixtures.


The tests include actual use of an induction hypothesis, nested $D$
formulas, transport through nested iterators, active witness erasure before
scope exit, ordinary elimination through a compound auxiliary, mismatched
seed/predecessor rejection, malformed constructors, Boolean indices,
unbound variables and non-object JSON envelopes. Both ordinary Python and
\texttt{python -O} are used. No trusted acceptance condition uses
\texttt{assert}.

The independent reviewer constructed further fixtures separately from the
main producers. The raw suite has five positives and seven adversaries,
each checked with and without optimization (24 expected outcomes). The
zero suite has two positives and four adversaries (12 expected outcomes).
All passed. These are finite implementation checks and independent replay,
not proof-assistant verification of the universal mathematical claims.
The general equivalence and both interpretation composites are proved
in the text.

\section{Limitations and further research questions}
\label{sec:questions}

The result removes an independent inference family and, in the main
presentation, all normalization from the trusted term specification. It
does not reduce arithmetic strength, discover a new primitive-recursive
function basis, or establish an absolute shortest proof checker. All the
inherited mechanisms have been credited: the scalar iterator, paired
binder, optional-antecedent clauses, selected introduction, Eval and
SeedEval, and the whole-domain arithmetic bridge.

The two computation patterns and identity fallback in $\rho$ are explicit.
The fixed reflexivity clause is explicit too. Moving a rule into an
unmeasured library or treating a normalizer as free would make an unfair
comparison; neither occurs here. The raw and zero checkers are independent
standalone acceptance programs, and their proof producers and translators
are untrusted supporting artifacts.

Several concrete questions remain:
\begin{enumerate}
\item \textbf{Formalize the fusion itself.} Add a deeply embedded
      three-family derivability relation to the pinned Lean development,
      and prove its equivalence with \texttt{VPrf} by the two local
      translations in Theorem~\ref{thm:raw-equivalence}. The raw route
      should be shorter because no normal-form conversion invariant is
      needed.
\item \textbf{Verify the inference-synthesizing checker.} Prove the
      correspondence between its negative free levels, positive bound
      indices, clause-count scope invariant and the embedded calculus.
      Its small source size is useful only when all those obligations
      remain visible.
\item \textbf{Reduce proof duplication.} The reverse Open translation
      may copy its major proof into the witness scope as a neutral
      auxiliary. A DAG certificate could share that proof, but the cost
      of reference validation and scope control must be measured.
\item \textbf{Replace the identity fallback economically.} Can the
      elimination simulation use a single always-active arithmetic pattern
      with no extra template and without enlarging either the checker or
      the arithmetic assumptions? A rule-count reduction alone would not
      answer this.
\item \textbf{Compare a fixed common certificate contract.} The compact
      codecs and the strict tree checkers sit on different tradeoffs.
      A useful next experiment would fix public syntax validation,
      certificate grammar, host semantics and measurement convention,
      then search for genuine dominance under that common contract.
\end{enumerate}

\section{Reproduction and attribution}

The archive includes standalone editable LaTeX, its compiled PDF, both
complete acceptance programs, matched controls, normalized sources,
proof transformations, examples and adversaries, measured results,
independent-review evidence and source provenance. Python reproduction
uses only the standard library:
\begin{lstlisting}[language={}]
python3 verify_manifest.py
python3 reproduce.py
python3 measure.py
python3 raw/checker.py < raw/examples/raw_zero_computation.json
python3 build_pdf.py
\end{lstlisting}
The PDF build requires the standard TeX packages named in the source and
runs with shell escape disabled. The main LaTeX source is flattened and
does not require the auxiliary section files. Rebuilding may change PDF
metadata; the manifest describes the delivered bytes.

The following appendices reproduce, with attribution, the five mathematical
foundation files from the pinned repository's
\texttt{Research/certified\_zero/reference/}. They provide the inherited
logical and arithmetic bridge in the delivered document. Their references
to an earlier ``new'' interface describe that historical interface, not
an additional novelty claim of this report. The present native primitives
are only those in Section~\ref{sec:raw}, or the explicitly stated zero
companion. Original copies and their hashes are retained in
\texttt{reference/foundations/}.

\clearpage
\appendix
\section{Complete raw acceptance checker}
The following is the entire measured 1,892-byte raw module. The public term syntax is unrestricted raw syntax; the two arithmetic cases are visible in the discharge branch.
\begin{lstlisting}[basicstyle=\ttfamily\fontsize{7.6}{8.7}\selectfont]
import json

def need(condition):
    if not condition:
        raise ValueError('invalid certificate')

def walk(a, f, d=0, form=None):
    if type(a) is int:
        need(not form)
        return f(a, d)
    need(type(a) is list)
    size = len(a)
    need(form is None or (size in (2, 1 if form else 3) and (form or size == 3 or a[0] == 0)))
    a = [walk(b, f, d + (2 if size == 1 else int(size == 3 and i == 2)), size == 1) for i, b in enumerate(a)]
    return a

def syntax(a, n, form=False):
    need(walk(a, lambda i, d: need(1 - n <= i <= d) or i, form=form) == a)

def sub(a, t, depth=0):
    return walk(a, lambda i, d: t if i == d + 1 else i - (i > d + 1), depth)

def prove(p, g):
    n = len(g)
    b = -n
    match p:
        case ['u', j, t, q]:
            need(type(j) is int and 0 <= j < n)
            syntax(t, n)
            u, B = g[j]
            if u:
                need(prove(q, g) == u)
            A = sub(B, t)
        case ['e', q, r, v]:
            D, = prove(q, g)
            B = sub(D, b, 1)
            local = g + [(sub(B, 0), B)]
            E = prove(r, local)
            A = prove(v, local)
            match (E, A):
                case [[s, t], [[0, a, h], z]] if s == a:
                    A = [z, t]
                case [[s, t], [[[0, c], a, h], z]] if s == [c, a, h]:
                    A = [z, sub(h, t)]
        case ['i', B0, t, r]:
            syntax(t, n)
            B = sub(B0, t, 1)
            u = sub(B, b if len(B0) == 2 and t == 0 else 0)
            need(prove(r, g + [([], u)]) == sub(B, [0, b]))
            A = [B0]
        case _:
            need(False)
    syntax(A, n, True)
    return A

def check(data):
    need(set(data) == {'goal', 'proof'})
    syntax(data['goal'], 1, True)
    need(prove(data['proof'], [([], [1, 1])]) == data['goal'])
if __name__ == '__main__':
    check(json.load(open(0)))
    print('ACCEPT')
\end{lstlisting}
\clearpage
\section{Complete zero-normal acceptance checker}
The following is the entire measured 1,858-byte companion module. Its traversal deletes zero-counter iterators, and its discharge branch uses the single SeedEval pattern.
\begin{lstlisting}[basicstyle=\ttfamily\fontsize{7.6}{8.7}\selectfont]
import json

def need(condition):
    if not condition:
        raise ValueError('invalid certificate')

def walk(a, f, d=0, form=None):
    if type(a) is int:
        need(not form)
        return f(a, d)
    need(type(a) is list)
    size = len(a)
    need(form is None or (size in (2, 1 if form else 3) and (form or size == 3 or a[0] == 0)))
    a = [walk(b, f, d + (2 if size == 1 else int(size == 3 and i == 2)), size == 1) for i, b in enumerate(a)]
    return a[1] if size == 3 and a[0] == 0 else a

def syntax(a, n, form=False):
    need(walk(a, lambda i, d: need(1 - n <= i <= d) or i, form=form) == a)

def sub(a, t, depth=0):
    return walk(a, lambda i, d: t if i == d + 1 else i - (i > d + 1), depth)

def prove(p, g):
    n = len(g)
    b = -n
    match p:
        case ['u', j, t, q]:
            need(type(j) is int and 0 <= j < n)
            syntax(t, n)
            u, B = g[j]
            if u:
                need(prove(q, g) == u)
            A = sub(B, t)
        case ['e', q, r, v]:
            D, = prove(q, g)
            B = sub(D, b, 1)
            local = g + [(sub(B, 0), B)]
            E = prove(r, local)
            A = prove(v, local)
            match (E, A):
                case [[s, t], [[[0, c], a, h], z]] if s == a:
                    A = [z, sub(h, t if c == 0 else [c, t, h])]
        case ['i', B0, t, r]:
            syntax(t, n)
            B = sub(B0, t, 1)
            u = sub(B, b if len(B0) == 2 and t == 0 else 0)
            need(prove(r, g + [([], u)]) == sub(B, [0, b]))
            A = [B0]
        case _:
            need(False)
    syntax(A, n, True)
    return A

def check(data):
    need(set(data) == {'goal', 'proof'})
    syntax(data['goal'], 1, True)
    need(prove(data['proof'], [([], [1, 1])]) == data['goal'])
if __name__ == '__main__':
    check(json.load(open(0)))
    print('ACCEPT')
\end{lstlisting}
\clearpage
% New mathematical result: complementary introduction choices.
% Requires amsmath/amssymb/amsthm and theorem, lemma, proposition environments.
% The inherited arithmetic-to-iPRA bridge must be supplied separately.

\section{Complementary introductions with one local assumption}
\label{sec:choice-introduction}

The old paired introduction retains both the zero instance and an
induction instance when its matrix is an equation. Both together are
unnecessary. A certificate may choose either one, provided the induction
choice is restricted to a literal equation. This changes neither the
formula language nor the meaning of a formula. The omitted combination
is recovered by a short nested introduction--elimination construction.

\subsection{Complete native syntax and rule inventory}

The scalar numerical language is
\[
 t::=x\mid0\mid St\mid I(n,a;u.h).
\]
Only $h$ binds the state variable $u$; $n$ and $a$ do not.
The update may contain fixed numerical parameters. The computation
axioms in every retained context are
\[
 I(0,a;u.h)=a,\qquad
 I(Sn,a;u.h)=h[I(n,a;u.h)/u].
 \tag{I0, IS}\label{eq:choice-computations}
\]
The argument also applies unchanged to the full fixed-description PR
language with its projection, composition, and recursion equations.
No numerical variable ranges over function descriptions, and no uniform
evaluator is postulated.

Native formulas and reusable context clauses are
\[
 A::=(s=t)\mid\mathsf D_{x,y}A,\qquad
 \langle A;y.B(y)\rangle,\quad y\notin\operatorname{FV}(A).
\]
Bound variables are renamed before capture-free substitution.
The two variables $x,y$ in a paired formula are distinct.
An ordinary assumption $F$ is the constant clause
$\langle H;v.F\rangle$, where $v$ is absent from $F$ and $H$ is the
fixed computation theorem $I(0,0;u.0)=0$. Thus ordinary hypothesis
use is a clause Use preceded by this computation proof.

There are computation axioms and exactly the following inference
patterns:
\[
\frac{\Gamma\vdash A}{\Gamma\vdash B(t)}
\quad\mathsf{Use}
\quad(\langle A;y.B(y)\rangle\in\Gamma),
\]
\[
\frac{\Gamma\vdash\mathsf D_{x,y}B(x,y)\qquad
      \Gamma,\langle B(a,0);y.B(a,y)\rangle\vdash F}
     {\Gamma\vdash F}
\quad\mathsf{DE},
\qquad
\frac{\Gamma\vdash s=t\qquad\Gamma\vdash F(s)}
     {\Gamma\vdash F(t)}
\quad\mathsf{Tr},
\]
\[
\frac{\Gamma,B(t,\eta_\varepsilon(z))\vdash B(t,Sz)}
     {\Gamma\vdash\mathsf D_{x,y}B(x,y)}
\quad\mathsf{DI}_{\varepsilon},
\qquad
\eta_0(z)=0,\quad\eta_1(z)=z.
\tag{Choice DI}\label{eq:choice-di}
\]
Here $\varepsilon$ is a certificate choice in $\{0,1\}$, and
$\varepsilon=1$ is permitted only when $B$ has a literal equation
root. An equation beneath a binder is not a literal equation.
Call $\varepsilon=0$ the \emph{case choice} and $\varepsilon=1$
the \emph{progress choice}. Both retain exactly one ordinary
formula assumption. Compound matrices have only the case choice.

In DI the eigenparameter $z$ is fresh for $\Gamma$, the witness
$t$, and every fixed parameter of $B$. In DE the eigenparameter
$a$ is fresh for $\Gamma$, the conclusion $F$, and every fixed
parameter of $B$. In transport, $F(v)$ is any fixed native
formula template. Displayed local assumptions and clauses are
discharged. There are no primitive reflexivity, coverage,
separation, ordinary logical, or additional induction rules.

We call this calculus $\mathcal C$. The original comparison
calculus $\mathcal N$ has identical syntax, computation,
Use, DE, and transport, but its introduction is
\[
\frac{\Gamma,B(t,0),[B(t,z)]_{\mathrm{at}}\vdash B(t,Sz)}
     {\Gamma\vdash\mathsf D_{x,y}B(x,y)}.
\tag{Old DI}\label{eq:choice-old-di}
\]
The bracket supplies $B(t,z)$ exactly at a literal equation.

\subsection{Exact equivalence and the unused-witness construction}

\begin{lemma}[Structural admissibility]
\label{lem:choice-structural}
Weakening, exchange, contraction, renaming, ordinary-assumption cut,
and capture-free numerical substitution are admissible in $\mathcal C$.
\end{lemma}
\begin{proof}
Induct on derivations. Rename local eigenparameters before inserting
proofs or terms. Numerical substitution preserves a formula's root
constructor, so an allowed progress choice remains allowed.
The selected choice does not change under substitution or cut.
For ordinary-assumption cut, replace each use of its constant clause
by the supplied proof. Computation schemas are closed under
capture-free numerical substitution. These facts verify every case.
\end{proof}

\begin{theorem}[Identity-language equivalence]
\label{thm:choice-equivalence}
For every retained clause context $\Gamma$ and native formula $F$,
\[
 \Gamma\vdash_{\mathcal C}F
 \quad\Longleftrightarrow\quad
 \Gamma\vdash_{\mathcal N}F.
\]
The transformations are effective on finite derivations and leave
all numerical and formula syntax unchanged.
\end{theorem}
\begin{proof}
Each new introduction is an instance of the old one after weakening
by the omitted local assumption. For the case choice at an equation,
add the induction instance. For the progress choice, add the zero
instance. Compound case introduction is already the old rule.
Thus every new derivation translates to an old one.

For the reverse direction, transform an old derivation recursively.
Only its atomic DI case needs a construction. Suppose its matrix is
the equation $E(x,y)$ and its transformed premise is
\[
 \pi:\quad\Gamma,E(t,0),E(t,z)\vdash E(t,Sz).
\]
The outer new inference will be case DI with witness $t$. It is
therefore enough to prove $E(t,Sz)$ under
$\Gamma,E(t,0)$.

Rename the eigenvariable in $\pi$ to a fresh $v$. Choose a new
variable $u$ absent from $E(t,y)$ and all retained data. Under
$\Gamma,E(t,0)$, progress DI with witness zero introduces
\[
             \mathsf D_{u,y}E(t,y)
\]
from exactly the renamed premise $\pi[v/z]$.
The bound coordinate $u$ is deliberately unused.

Apply DE to this paired formula. Because $u$ is unused, the opened
clause is exactly
\[
               \langle E(t,0);y.E(t,y)\rangle;
\]
it contains no new existential witness. The retained ordinary
assumption $E(t,0)$ supplies its antecedent. Use at $Sz$ gives
$E(t,Sz)$, and DE closes. Finally the outer case DI introduces
the original $\mathsf D_{x,y}E(x,y)$ with witness $t$.
The original premise proof $\pi$ has been inserted once.

An old compound introduction already is the new case choice.
All other rule cases are unchanged. This completes the induction
on the original derivation.
\end{proof}

The unused inner coordinate is essential. Introducing
$\mathsf D_{x,y}E(x,y)$ internally and then eliminating it would
give an arbitrary fresh witness, not necessarily the original
$t$. The construction instead keeps $t$ as a fixed parameter
of $E(t,y)$ and binds only the unused coordinate $u$.

\begin{proposition}[Linear bound on inference nodes]
\label{prop:choice-node-bound}
Count each computation axiom, Use, DI, DE, and transport as one
inference node, expanding ordinary-assumption use as a computation
of $H$ followed by Use. If an old derivation has $N$ nodes and
$r$ atomic DI nodes, the displayed transformation has exactly
\[
                   N+5r\leq 6N
\]
nodes.
\end{proposition}
\begin{proof}
Each atomic old DI node is replaced by six nodes: outer case DI,
inner progress DI, DE, Use of the opened clause, Use of the
retained base assumption, and the computation of $H$ needed for
that last ordinary-assumption use. Its old premise subtree
occurs once. All other inference nodes are copied once.
\end{proof}

This counts inference-tree nodes. It is not a bound on serialized
JSON bytes, formula sizes, checker runtime, or the number of
consumed packet cells when hypothesis closure is implicit.
In the formula-context Q/D presentation, the corresponding
replacement has outer case DI, inner progress DI, DE, QE,
and two Hyp nodes, giving the same six-node count.

\subsection{Induction and numerical cases separate cleanly}

\begin{lemma}[Contextual equation induction]
\label{lem:choice-ei}
The progress choice, DE, and Use derive
\[
\frac{\Gamma\vdash E(0)\qquad
      \Gamma,E(z)\vdash E(Sz)}
     {\Gamma\vdash E(t)}
\tag{EI}\label{eq:choice-ei}
\]
for every equation $E$, with the usual freshness condition.
\end{lemma}
\begin{proof}
Choose an unused witness coordinate $u$. Progress DI introduces
$\mathsf D_{u,y}E(y)$ from the given step. DE opens the clause
$\langle E(0);y.E(y)\rangle$. Supply its antecedent using the given
base proof and use the clause at $t$. Close DE.
\end{proof}

\begin{lemma}[Cases, equality, and binder congruence]
\label{lem:choice-preinduction}
The case choice derives, for arbitrary native formulas,
\[
\frac{\Gamma\vdash F(0)\qquad\Gamma\vdash F(Sz)}
     {\Gamma\vdash F(t)}
\tag{NC}\label{eq:choice-nc}
\]
with the usual freshness condition. Equality reflexivity,
symmetry, transitivity, numerical congruence, and uniform
matrix interderivability under $\mathsf D$ are derivable
without using the progress choice.
\end{lemma}
\begin{proof}
For cases, use case DI with an unused witness coordinate and
matrix $F(y)$, ignoring its retained base in the successor
premise. DE opens $\langle F(0);y.F(y)\rangle$; the given zero
proof and Use at $t$ finish.

For reflexivity put $r=I(0,t;u.0)$. Use the computation equation
$r=t$ as both transport premises with motive $v=t$ to obtain
$t=t$. The usual explicit equality motives give symmetry,
transitivity, and numerical congruence.

For binder congruence, open the old paired formula at a fresh
witness $a$ and introduce the new one by case DI at $a$.
Convert its retained base to the old base, use the opened
old clause at the successor, and convert back. Ordinary cut
closes the construction, and the reverse conversion gives
the converse. No induction is used.
\end{proof}

The primitive progress choice is still restricted to equations.
Allowing it at arbitrary compound matrices would, by the
argument for \eqref{eq:choice-ei}, supply induction for
arbitrary formulas and change the intended theory.

\subsection{The entire ordinary logical interface}
\label{subsec:choice-logic}

Here is a predecessor-free reconstruction of the ordinary logical
operations. It makes explicit that the case choice supplies the
logical interface independently of progress induction.

Let $C(v,a,b):=I(v,a;u.b)$ with $u$ unused. Computation gives
$C(0,a,b)=a$ and $C(Sv,a,b)=b$. Define the gate on fixed finite
formula syntax by
\[
 G_v(s=t):=(C(v,0,s)=C(v,0,t)),\qquad
 G_v(\mathsf D_{x,y}B):=\mathsf D_{x,y}G_vB,
\]
renaming bound variables away from the control. Equality,
binder congruence, and external structural induction give
\[
 \vdash G_0A,\qquad G_{Sv}A\dashv\vdash A,\qquad
 A\dashv\vdash\mathsf D_{u,y}G_yA
 \quad(u,y\notin\operatorname{FV}(A)).
\tag{Gate and padding}\label{eq:choice-gate}
\]
For padding, case DI proves the forward direction. DE, the
zero-gate theorem, and Use at $1$ prove the reverse direction.

Define a selector between two arbitrary formulas by
\[
 M_v(s=t,u=w):=(C(v,s,u)=C(v,t,w))
\]
and, with aligned fresh paired variables, by
\[
\begin{aligned}
 M_v(\mathsf D_{x,y}A,\mathsf D_{x,y}B)
   &:=\mathsf D_{x,y}M_v(A,B),\\
 M_v(E,\mathsf D_{x,y}B)
   &:=\mathsf D_{x,y}M_v(G_yE,B),\\
 M_v(\mathsf D_{x,y}A,E)
   &:=\mathsf D_{x,y}M_v(A,G_yE),
\end{aligned}
\]
where $E$ is an equation. Every call decreases maximum input
binder depth. Padding and congruence prove
$M_0(A,B)\dashv\vdash A$ and
$M_{Sv}(A,B)\dashv\vdash B$.

Put $\delta(v)=C(v,1,0)$ and $p(z)=I(z,0;u.\delta(u))$.
Only the fixed computations $p(1)=1$ and $p(2)=0$ are needed.
The full logical macro inventory is
\[
\begin{aligned}
 A\Rightarrow_*B&:=\mathsf D_{u,z}M_z(A,B),\\
 A\wedge_*B&:=\mathsf D_{u,z}G_zM_{p(z)}(A,B),\\
 \exists_*x\,A(x)&:=\mathsf D_{x,y}G_yA(x),\\
 A\vee_*B&:=\exists_*x\,M_x(A,B),\\
 \forall_*x\,A(x)&:=
        A(0)\wedge_*\mathsf D_{u,z}G_zA(z),\\
 \bot_*&:=(0=1).
\end{aligned}
\tag{Logical macros}\label{eq:choice-macros}
\]
Newly bound variables are fresh; $u$ is unused.

Implication introduction obtains its assumption from the zero
selector and supplies the consequent at successors, then uses
case DI. Elimination opens its clause and uses it at $1$.
Existential introduction uses its witness and the successor
gate; elimination opens at a fresh witness, supplies the
zero-gate theorem, and uses the clause at $1$.

Disjunction introduction uses witnesses $0$ and $1$.
For elimination, discharge the two branch proofs as
$A\Rightarrow_*F$ and $B\Rightarrow_*F$. NC applied to
$M_z(A,B)\Rightarrow_*F$ gives this implication at every
control: the two endpoint cases follow from the branch
implications. Existential elimination provides $M_a(A,B)$
and the implication gives $F$.

Given $A$ and $B$, selector endpoints and NC give $M_v(A,B)$
at an arbitrary control. Substitute $v=p(Sz)$ and use the
successor gate for conjunction introduction. For projections,
open its zero-gated clause and specialize at $1$ and $2$;
the two fixed computations give $B$ and $A$, respectively.

These rules prove existential numerical coverage:
\[
     (x=0)\vee_*\exists_*v\,(x=Sv).
\tag{Coverage}\label{eq:choice-coverage}
\]
Use NC. Zero takes the left injection of reflexivity, and a
successor takes the right injection with the evident witness.
Coverage gives equality cases even when the retained context
contains $x$; its elimination adds $x=0$ or $x=Sv$ rather
than silently changing that retained context.

Universal introduction follows from a fresh-parameter proof
by substituting zero and successors and forming the two
conjuncts. For elimination, take its $A(0)$ conjunct and
open the paired conjunct. Its base is a zero-gate theorem.
Coverage of the target chooses the zero instance from
$A(0)$ or a successor instance from the opened clause
and gate removal.

Finally, under $0=1$, computation and transport give
$s=C(0,s,t)=C(1,s,t)=t$ for arbitrary numerical terms.
Every equation follows, and external induction on formula
syntax with case DI gives every native formula. This
supplies falsity elimination. No excluded middle for
quantified formulas has entered the construction.

\subsection{Whole-domain realizations and both formula composites}

For a standard first-order reading, let $T_{\rm eq}$ have
intuitionistic equality logic, the same numerical signature
and computation equations, numerical coverage, separation
$(0=1)\to\bot$, and contextual equation induction.
Its interpretation of the native syntax is
\[
 (s=t)^\circ=(s=t),\qquad
 (\mathsf D_{x,y}B)^\circ
   =\exists x\bigl(B^\circ(x,0)\to\forall y B^\circ(x,y)\bigr),
\]
\[
 \langle A;y.B\rangle^\circ
       =A^\circ\to\forall y B^\circ.
\]
The domain is the entire numerical domain and interpreted
equality is actual equality.

Case DI is sound in this realization by numerical coverage.
Progress DI is sound by contextual equation induction under
the conditional base. Use, DE, transport, and computation
have their displayed ordinary meanings. Conversely, the
native macros derive every ordinary logical rule,
\eqref{eq:choice-ei} derives equation induction, and
\eqref{eq:choice-coverage} derives coverage. Separation
translates to the identity implication
$(0=1)\Rightarrow_*(0=1)$.

For the essential paired formula composite,
\[
 \mathsf D_{x,y}B
 \dashv\vdash
 \exists_*x\bigl(B(x,0)\Rightarrow_*\forall_*y\,B(x,y)\bigr).
\tag{Paired composite}\label{eq:choice-paired-composite}
\]
Forwards, DE opens its clause at a fresh $a$; under $B(a,0)$
that clause supplies $B(a,y)$ at a fresh arbitrary $y$.
Introduce the universal, implication, and existential
macros and close DE. Backwards, eliminate the existential
at a fresh $a$. Introduce the paired formula with witness
$a$ by case DI: the retained base, supplied implication,
and universal elimination at $Sz$ give its successor premise.

The ordinary interpretations of the six logical macros
have their ordinary round trips. For conjunction, its true
base and the instances $1,2$ recover both operands;
the converse uses selector cases. The universal has all
successor instances in its paired conjunct and zero in
its first conjunct; coverage gives exactly the ordinary
universal. Implication, disjunction, and existential
quantification use their selector or gate endpoints.
Falsity uses the explicit ordinary separation bridge.
External induction on finite formulas, together with
\eqref{eq:choice-paired-composite}, gives both formula
composites. Clauses correspond in both directions to
$A\Rightarrow_*\forall_*yB$ by the derived logical rules.

The numerical composites are literal identities with
internally provable graph $x=y$. No standard-numeral
restriction, retract, quotient, or equality reflection
assumption is being inserted at this stage. The separately
supplied arithmetic bridge identifies $T_{\rm eq}$ with
full iPRA; the present change preserves that bridge
unchanged because Theorem~\ref{thm:choice-equivalence}
preserves every native sequent.

\subsection{What the theorem does and does not claim}

The new inference interface retains one ordinary assumption
at every paired introduction. It preserves the entire
consequence relation, not just equations or closed
theorems. Consequently the established full-iPRA
interpretations and equiconsistency transfer unchanged.

This result is different from retaining only progress
introduction and strengthening closure elimination.
Here the original elimination rule is retained, and
the case and progress introductions coexist. It is also
different from allowing arbitrary compound progress,
which would supply stronger induction.

A complete-checker comparison must still count the
certificate choice, its equation-root restriction,
substitution, context handling, computation, and the
public input/output interface. The mathematical theorem
alone is not a byte-count or minimality theorem.
The inference-node bound is distinct from source-size
measurements and from serialized certificate size.
The proof is a written universal proof; finite regression
tests do not constitute a Lean/Rocq formalization of it.

\section{Arithmetic strength and both ordinary interpretation composites}
\label{sec:ordinary-bridge}

We spell out the arithmetic bridge to avoid identifying standard numerical
agreement with bi-interpretability. First elaborate every fixed iterator
into an ordinary function symbol. At stage zero the symbols are $0,S$.
Given a fixed earlier-stage term $h(u,\vec p)$, introduce
$i_h(n,a,\vec p)$ with the two equations
\[
i_h(0,a,\vec p)=a,\qquad
i_h(Sn,a,\vec p)=h(i_h(n,a,\vec p),\vec p).
\]
Only finite descriptions occur; the domain contains no code for a
universal evaluator. Formal parameters are distinct when a description
is declared, while actual arguments can coincide.

Let $T_I$ be ordinary intuitionistic first-order logic with actual
equality and genuine falsity, these staged equations, contextual
induction for literal equations, and
\[
\forall x\,(x=0\vee\exists y\,x=Sy),
\qquad (0=S0)\to\bot.
\]
The second sentence is the falsity bridge. It must be stated: positive
numerical equations by themselves admit a singleton structure. Native
explosion from $0=S0$ and a translation of ordinary falsity to that
equation do not authorize silently dropping the bridge from $T_I$.

The native-to-ordinary translation is
\[
(s=t)^\circ=(s=t),\qquad
(\mathsf D_{x,y}B)^\circ=
\exists x\bigl(B^\circ(x,0)\to\forall y B^\circ(x,y)\bigr).
\]
A clause $\langle A;y.B\rangle$ becomes $A^\circ\to\forall yB^\circ$.
For the auxiliary closure used by the compact checkers,
$(\mathsf Q_yB)^\circ=B^\circ(0)\to\forall yB^\circ$.
The interpretation domain is all elements, defined by $x=x$;
equality is actual equality.

Case DI is sound because numerical coverage combines its retained base
with its successor proof. Progress DI is sound by literal-equation
induction under a temporarily assumed base. Use, DE and transport are
ordinary logical inferences under the translations. Conversely the
derived logical operations, equation induction, cases and explosion in
Section~\ref{sec:choice-introduction} recover the rules and axioms of
$T_I$. An auxiliary closure is interderivable with
$\mathsf D_{v,y}B(y)$ for unused $v$: in one direction use case DI with
the closure and its base; in the other eliminate the paired formula
and return its supplied closure. Thus formulas and context clauses have
the stated meanings in both presentations.

\subsection{The arithmetic dependency order}

The scalar update has no recursion-counter argument. We do not conceal
one in an encoding. The inherited internal proofs, reproduced in the
appendices, establish the required arithmetic in this order:
\begin{enumerate}
\item Using equation induction and numerical cases, construct a fixed
scalar predecessor term and prove $P0=0$ and $P(Sn)=n$.
Existential coverage then gives $x=0\vee x=S(Px)$: in the successor
case $x=Sy$, apply $P$ to obtain $Px=y$ and substitute.
\item Construct truncated subtraction and a numerical characteristic
of actual equality. Prove its reflection on the entire original domain.
\item Compile each fixed quantifier-free formula into its characteristic
equation and recover every contextual quantifier-free induction instance.
\item With that induction available, construct pairing and projections
by a bounded square-root search and compile every fixed general
primitive-recursive description into a scalar iterator term.
\end{enumerate}
No pairing or faithful equality detector is assumed to obtain
predecessor. No appeal to standard-model correctness replaces an
internal equality. Every induction fixes its passive parameters;
diagonal substitution occurs only into a previously proved identity.

For a target recursion
\[
f(0,\vec p)=g(\vec p),\qquad
f(Sn,\vec p)=h(n,f(n,\vec p),\vec p),
\]
let $\pi,L,H$ be the proved pairing and projections. The compiler uses
\[
V(n,\vec p)=I\bigl(n,\pi(0,g^*(\vec p));
u.\pi(SL(u),h^*(L(u),H(u),\vec p))\bigr),
\qquad f^*(n,\vec p)=H(V(n,\vec p)).
\]
Equation induction proves $L(V(n,\vec p))=n$; the two projection
identities then prove the target defining equations. The construction
is external recursion on each finite PR description, not an internal
quantification over arbitrary functions.

\begin{theorem}[Whole-domain bi-interpretability and equiconsistency]
The ordinary realization $T_I$ and full first-order iPRA are
bi-interpretable. In both theories $x=y$ defines the isomorphism from
the original domain to its double interpretation. The native Choice-DI
calculus is therefore a faithful proof presentation of a theory
bi-interpretable and equiconsistent with iPRA.
\end{theorem}
\begin{proof}
Interpret an iterator in iPRA by primitive recursion that ignores its
counter argument in the update. iPRA proves its equations, literal
equation induction, coverage and separation. In the reverse direction
use the scalar compiler. The predecessor and equality-reflection proofs
supply the successor axioms and every quantifier-free induction instance.
All ordinary logical rules are preserved.

For two iterators with the same seed and provably equal updates,
equation induction on the equality of their values proves uniqueness at
every counter, with all parameters fixed. The corresponding argument
applies to ordinary primitive recursion. Induct externally on a finite
description: its double translation satisfies its original defining
equations, so internal uniqueness identifies the two functions. This
also proves coherence when substitution changes the staged syntactic
description of an update.

The domain and equality are unchanged; hence the identity graph is
internally total, bijective and preserves every interpreted operation.
The displayed ordinary meanings of the logical macros and the paired
binder give both formula composites by external induction on finite
formula syntax. Clauses are similarly equivalent to their conditional
universal meanings. The explicit falsity bridge preserves contradiction
in both directions. Effective translations of finite proofs give
equiconsistency, without a claim that either theory proves its own
consistency.
\end{proof}

% Adapted with attribution from the pinned unparenthesized one-hole archive.
\section{Internal scalar predecessor bootstrap}
\label{sec:hybrid-bootstrap}

This scalar bootstrap is inherited from the supplied scalar development~\cite{erasure} and is reproduced to keep the scalar checker self-contained. It is not a new contribution of this paper. The proof uses a factor map fixing odd states and a saturated-iteration construction of predecessor.

The available arithmetic terms are $0$, $S$, and
\[
 I(n,a;u.h),\qquad
 I(0,a;u.h)=a,\qquad
 I(Sn,a;u.h)=h[I(n,a;u.h)/u].
\]
The update binds only $u$.  All other variables in it are fixed numerical
parameters.  We use equality reasoning, the already reconstructed
zero/successor cases, and induction on a single literal equation.
Every induction below keeps its other parameters fixed.  Predecessor,
successor injectivity, pairing, and induction on arbitrary
quantifier-free formulas are not assumptions of the argument.

\subsection{Elementary terms and two finite cycles}

Write $1=S0$ and use larger numerals as successor abbreviations.  Define
\[
 C(n,a,b)=I(n,a;u.b),\qquad
 Z(n)=C(n,1,0),\qquad
 a+b=I(b,a;u.Su).
\]
The update variable of $C$ is absent from its body.  Consequently
\[
 C(0,a,b)=a,\qquad C(Sn,a,b)=b.
\]
Numerical cases show that $Z(n)$ is $0$ or $1$.
The usual commutative-monoid equations for addition follow by literal
equation induction.  For example, $0+b=b$ and $Sa+b=S(a+b)$ are
proved by induction on $b$; these yield commutativity by a further
equation induction, and associativity follows by induction on the last
summand.  We use these equations to reorder finite sums.  The notation
$2a$ means $a+a$ and $3a$ means $a+a+a$; no multiplication operation
is needed for this bootstrap.

Put
\[
 N(n)=I(n,0;u.Z(u)),\qquad
 j(u)=Z(u)+2N(u),\qquad
 M(n)=I(n,0;u.j(u)).
\]
The range of $N$ is contained in $\{0,1\}$: its zero value is $0$,
and its successor value is a value of $Z$.  This uses numerical cases
on its counter, not induction on a disjunction.  Hence
\begin{equation}
 N(n)+Z(N(n))=1,\qquad
 Z(Z(N(n)))=N(n),\qquad N(n+2)=N(n).
 \label{eq:hybrid-parity}
\end{equation}
Here the last equation follows from two successor computations and
the preceding equation.  Literal-equation inductions give
\begin{equation}
 N(2n)=0,\qquad N(2n+1)=1,\qquad N(3n)=N(n).
 \label{eq:hybrid-parity-scaling}
\end{equation}
For the first equation the successor step adds two to the argument of
$N$; the second follows by one computation.  In the third, increasing
$n$ adds three to the left argument and therefore applies $Z$ once to
both sides of the induction equation.

Finite computations now give
\[
 j(0)=1,\qquad j(1)=2,\qquad j(2)=0.
\]
The proof below needs just these three values and
$M(Sn)=j(M(n))$; it does not assume a division algorithm or a theorem
characterizing $M$ as remainder modulo three.

\subsection{A ramp with a global scaling identity}

Define
\[
 G(n)=I(n,0;u.u+M(u)+1).
\]
The central invariant is
\begin{equation}
 M(G(n))=N(n).
 \label{eq:hybrid-ramp-invariant}
\end{equation}
Its zero case is immediate.  For its successor step, retain the
single induction equation and split the Boolean value $N(n)$.
If it is $0$, the next $G$ state is $G(n)+1$; its $M$ value is
$j(0)=1$, equal to $N(Sn)$.  If it is $1$, the next $G$ state
is $G(n)+2$; two computations of $M$ change its value from $1$
to $2$ and then to $0$, again equal to $N(Sn)$.  This proves
the invariant by literal-equation induction.

It follows that
\[
 G(Sn)=G(n)+N(n)+1,\qquad G(n+2)=G(n)+3.
\]
Equation induction therefore yields
\begin{equation}
 G(2n)=3n.
 \label{eq:hybrid-ramp-double}
\end{equation}
A separate literal-equation induction proves
\begin{equation}
 G(3n)=3G(n)+N(n).
 \label{eq:hybrid-ramp-triple}
\end{equation}
For completeness, its successor calculation is
\begin{align*}
 G(3Sn)
 &=G(3n)+N(3n)+4
   =3G(n)+2N(n)+4,\\
 3G(Sn)+N(Sn)
 &=3G(n)+3N(n)+3+Z(N(n)).
\end{align*}
The expressions agree in both Boolean cases for $N(n)$.
The base is $0=0$.  In particular,
\eqref{eq:hybrid-ramp-triple} was proved internally; it was not
inferred from an intended floor-function interpretation of $G$.

Set
\[
 H(u)=C(N(u),G(u),u).
\]
The crucial choice is to retain $u$ when its parity is $1$.
The preceding equations give
\begin{equation}
 H(2u)=3u,\qquad H(1)=1,\qquad H(3u)=3H(u).
 \label{eq:hybrid-factor-map}
\end{equation}
To check the last identity, use $N(3u)=N(u)$ and split the Boolean
$N(u)$.  At $0$, both sides are $3G(u)$ by
\eqref{eq:hybrid-ramp-triple}; at $1$, both are $3u$.
This argument does not require that an arbitrary number has already
been represented as $2v$ or $2v+1$.

\subsection{A comparator from scalar iteration}

For every fixed update $h$, equation induction on $n$, keeping $m,a$
and all parameters fixed, proves the semigroup identity
\[
 I(m+n,a;h)=I(n,I(m,a;h);h).
\]
The base is the zero computation and the step applies $h$ to the
induction equation. This identity will also be used in the later
pairing construction.


For a fixed update $h$, equation induction on $y$, with its initializer
$a$ and all parameters of $h$ fixed, proves the rotation identity
\begin{equation}
 I(Sy,a;h)=I(y,h(a);h).
 \label{eq:hybrid-rotation}
\end{equation}
Both sides have base $h(a)$, and their successor equations apply the
same $h$ to the previous equality.  Here and below $I(y,a;h)$ is
an abbreviation for one fixed finite scalar term, not quantification
over functions.

By \eqref{eq:hybrid-factor-map}, another equation induction gives
\begin{equation}
 I(y,3a;H)=3I(y,a;H).
 \label{eq:hybrid-iterator-scaling}
\end{equation}
Its base is $3a=3a$.  In the step the fixed-parameter induction
equation rewrites the new left side as
\[
 H(3I(y,a;H))=3H(I(y,a;H)).
\]
The initializer parameter stays $a$ throughout this induction.
Equation induction also gives $I(y,1;H)=1$, since $H(1)=1$.

Define
\[
 E(x)=I(x,1;u.2u),\qquad
 L(x,y)=N\bigl(I(y,E(x);H)\bigr).
\]
Then the following entire recursion specification is provable:
\begin{equation}
 L(0,y)=1,\qquad L(Sx,0)=0,\qquad L(Sx,Sy)=L(x,y).
 \label{eq:hybrid-comparator}
\end{equation}
The first equation uses $E(0)=1$ and the fixed-point identity for
$H$.  The second uses $E(Sx)=2E(x)$ and $N(2a)=0$.
For the third,
\begin{align*}
 L(Sx,Sy)
 &=N\bigl(I(Sy,2E(x);H)\bigr)\\
 &=N\bigl(I(y,H(2E(x));H)\bigr)\\
 &=N\bigl(I(y,3E(x);H)\bigr)\\
 &=N\bigl(3I(y,E(x);H)\bigr)\\
 &=N\bigl(I(y,E(x);H)\bigr)=L(x,y).
\end{align*}
The successive steps use rotation, $H(2a)=3a$, iterator scaling,
and $N(3a)=N(a)$.  Thus $L$ has the comparison equations at all
arguments internally.  Its standard-model meaning as the indicator
of $x\le y$ is motivation, not a premise of the following argument.

\subsection{Saturated iteration gives predecessor}

We now need only the three equations
\eqref{eq:hybrid-comparator}.  Define
\[
 B_m(u)=C(L(Su,m),u,Su),\qquad
 W(t,m)=I(t,0;u.B_m(u)).
\]
The subscript in $B_m$ records a fixed parameter; $B_m$ is an open
term abbreviation, not a function variable.

Numerical cases on $e$ prove the selector equation
\[
 C(e,Sa,Sb)=S C(e,a,b).
\]
Equivalently one may use equation induction on $e$, whose induction
hypothesis is unused.  Therefore
\begin{align}
 B_{Sm}(Su)
 &=C(L(SSu,Sm),Su,SSu)\notag\\
 &=C(L(Su,m),Su,SSu)=S B_m(u),
 \label{eq:hybrid-conjugacy}\\
 B_{Sm}(0)&=C(L(S0,Sm),0,S0)=C(L(0,m),0,S0)=S0.
 \label{eq:hybrid-conjugacy-base}
\end{align}

Literal-equation induction on $t$, with $m$ fixed, proves
\begin{equation}
 W(St,Sm)=S W(t,m).
 \label{eq:hybrid-saturated-shift}
\end{equation}
The base is \eqref{eq:hybrid-conjugacy-base}.  The successor step is
\begin{align*}
 W(SSt,Sm)
 &=B_{Sm}(W(St,Sm))\\
 &=B_{Sm}(S W(t,m))\\
 &=S B_m(W(t,m))=S W(St,m).
\end{align*}
Only the second line uses the current induction hypothesis, at the
same fixed $m$.

Two subsequent equation inductions give
\begin{equation}
 L(Sn,n)=0,\qquad W(n,n)=n.
 \label{eq:hybrid-diagonals}
\end{equation}
The first uses the comparator recurrence.  The second uses
\eqref{eq:hybrid-saturated-shift} in its successor step.  The
substitution $t=m=n$ is made into that already proved universally
parameterized equation; no parameter of an active induction hypothesis
is changed.  Consequently
\begin{equation}
 W(Sn,n)=B_n(n)=C(L(Sn,n),n,Sn)=n.
 \label{eq:hybrid-extra-step}
\end{equation}

Finally define the finite scalar term
\begin{equation}
 P(m)=I\bigl(m,0;u.C(L(SSu,m),u,Su)\bigr).
 \label{eq:hybrid-predecessor-term}
\end{equation}
Its zero equation is immediate.  To justify its successor equation
without assuming a separate extensionality rule for iterator bodies,
put
\[
 V(t,m)=I\bigl(t,0;u.C(L(SSu,Sm),u,Su)\bigr).
\]
The comparator recurrence identifies this update pointwise with
$B_m(u)$.  Equation induction on $t$, with $m$ fixed, gives
$V(t,m)=W(t,m)$: their base values are $0$, and the successor step
uses the pointwise update equation and then the induction equation.
Now \eqref{eq:hybrid-extra-step} yields
\[
 \boxed{\qquad P(0)=0,\qquad P(Sm)=V(Sm,m)=W(Sm,m)=m.\qquad}
\]

\subsection{Dependency boundary}

All terms in this section are built from $0,S,I$ by finite substitution.
All inductions have a single equation as their invariant.  Numerical
cases and equality transport are the already reconstructed logical
resources.  Successor injectivity may now be obtained by applying $P$
to an assumed equation $Sx=Sy$.  Equality tests, reflection of
quantifier-free formulas, pairing, and the recovery of general
primitive recursion can then be developed in the established order.
They were not used in order to manufacture predecessor.

This argument is a written internal derivation in the stated theory.
The finite implementation checks do not constitute a proof-assistant
formalization of this entire metatheorem.

% Standalone section.  Requires only amsmath and amssymb in the parent document.
% No theorem environments, macros, or bibliography keys are assumed.
% All arithmetic claims in this section are proved here.

\section{From literal-equation induction to full intuitionistic PRA}
\label{sec:eic-completion}\label{app:arithmetic}

This section supplies an unconditional completion theorem. It is adapted
with attribution from the current forward-left draft~\cite{forward}. In the
scalar application used here, every function occurring in the proof is
already a scalar term: predecessor is supplied by
Appendix~\ref{sec:hybrid-bootstrap}, the selector is
$C(n,a,b)=I(n,a;u.b)$, addition is $x+y=I(y,x;u.Su)$, and
truncated subtraction is $d(x,y)=I(y,x;u.Pu)$. All other functions below
are finite compositions of these. Thus this proof does not presuppose
that arbitrary primitive recursion has already been compiled. The QFI
conclusion first applies to the scalar language; the general PR compiler
comes afterward.  Literal-equation
induction, together with the primitive-recursive defining equations,
predecessor coverage, and separation of $0$ and $1$, proves induction for every
quantifier-free formula.  In particular, the proof does not assume a faithful
numerical equality test, stability or decidability of equality, successor
uniqueness, or any commutation law for an iterator.  We construct the equality
test and prove its reflection property for arbitrary elements of the
numerical domain.

\subsection{The exact hypotheses}

Work in ordinary intuitionistic first-order logic with equality and genuine
falsity $\bot$.  The language has one numerical sort and a symbol for each
fixed, finite primitive-recursive description.  Distinct descriptions are
not identified merely because they compute the same standard function.
Let $\mathsf{PRE}$ contain the defining equations for zero, successor,
projections, composition, and primitive recursion.  In particular, choose
descriptions $P$ and $C$ with equations
\begin{align}
 P(0)&=0, & P(Sx)&=x, \label{eq:eic-predecessor}\\
 C(0,a,b)&=a, & C(Sx,a,b)&=b. \label{eq:eic-selector}
\end{align}
Write $1=S0$.  In addition to $\mathsf{PRE}$, assume
\begin{align}
 \mathsf{Cov}:&\qquad
       \forall x\,\bigl(x=0\lor x=S(Px)\bigr),
       \label{eq:eic-coverage}\\
 \mathsf{Sep}:&\qquad 0=1\longrightarrow\bot,
       \label{eq:eic-separation}
\end{align}
and induction for each \emph{literal equation} $E(n,\vec p)$:
\begin{equation}
 \forall\vec p\,
 \left[
   \left(E(0,\vec p)\land
       \forall n\bigl(E(n,\vec p)\longrightarrow E(Sn,\vec p)\bigr)
   \right)
   \longrightarrow\forall n\,E(n,\vec p)
 \right].
 \label{eq:eic-induction}
\end{equation}
Here the schematic expression $E$ is a single equality of numerical
terms; the displayed universal quantifiers range only over numbers.
Every description and every instance of the schema is finite external
syntax.

The axiom formulation \eqref{eq:eic-induction} and the usual contextual
induction rule are interchangeable.  For example, if
\[
 \Gamma\vdash E(0,\vec p),
 \qquad
 \Gamma,E(z,\vec p)\vdash E(Sz,\vec p),
\]
where $z$ is absent from $\Gamma$ and the passive parameters, discharge the
second premise's equation and universally introduce $z$.  The axiom then
gives $\Gamma\vdash\forall n E(n,\vec p)$, which can be instantiated at any
target term.  Conversely, assume the antecedent of the axiom and apply the
contextual rule with fresh variables.  This uses the ordinary intuitionistic
logical rules and permits arbitrary, including quantified, formulas in
$\Gamma$.

\subsection{Constructor separation and two elementary tests}

Define the sign description and the countdown description by
\begin{equation}
 \operatorname{sg}(u)=C(u,0,1),
 \qquad
 K(u)=C(u,1,Pu).
 \label{eq:eic-sign-countdown}
\end{equation}
Their defining equations imply
\begin{align}
 \operatorname{sg}(0)&=0,
 &\operatorname{sg}(Sx)&=1,\label{eq:eic-sign-computation}\\
 K(0)&=1,
 &K(Sx)&=x.\label{eq:eic-countdown-computation}
\end{align}

\paragraph{Constructor laws.}
Successor is injective: applying $P$ to $Sx=Sy$ gives $x=y$ by
\eqref{eq:eic-predecessor}.  Successor never equals zero.  Indeed,
$Sx=0$ gives $1=0$ by applying $\operatorname{sg}$ and using
\eqref{eq:eic-sign-computation}; symmetry and
\eqref{eq:eic-separation} give $\bot$.  Consequently $1\ne0$, where
$A\ne B$ abbreviates $A=B\to\bot$.  None of these arguments uses
induction.

\paragraph{The sign test.}
For every $u$,
\begin{align}
 &\operatorname{sg}(u)=0\lor\operatorname{sg}(u)=1,
       \label{eq:eic-sign-range}\\
 &\operatorname{sg}(u)=0\quad\longleftrightarrow\quad u=0,
       \label{eq:eic-sign-zero}\\
 &\operatorname{sg}(u)=1\quad\longleftrightarrow\quad u\ne0.
       \label{eq:eic-sign-one}
\end{align}
To prove the range assertion, split \eqref{eq:eic-coverage}.  The branch
$u=0$ gives sign $0$, and the branch $u=S(Pu)$ gives sign $1$.
The implication $u=0\to\operatorname{sg}(u)=0$ is computation.
For its converse, the successor branch of coverage would make
$\operatorname{sg}(u)=0$ imply $1=0$, a contradiction; in that branch
explosion gives $u=0$ as well.  Similarly, if
$\operatorname{sg}(u)=1$ and $u=0$, computation gives $0=1$.
Conversely, $u\ne0$ excludes the zero branch of coverage, and the
successor branch computes the sign as $1$.

Coverage and \eqref{eq:eic-selector} also give, for all $u,v$,
\begin{equation}
 C(u,v,v)=v.
 \label{eq:eic-selector-idempotence}
\end{equation}
Both coverage branches compute to $v$, so this identity again needs no
induction.

\paragraph{The countdown test and the notation $W$.}
For every $u$,
\begin{equation}
 K(u)=0\quad\longleftrightarrow\quad u=1.
 \label{eq:eic-countdown-zero}
\end{equation}
If $u=1=S0$, \eqref{eq:eic-countdown-computation} gives $K(u)=0$.
Conversely, assume $K(u)=0$ and split coverage.  If $u=0$, then
$K(u)=1$, contradicting separation, and explosion gives $u=1$.
If $u=S(Pu)$, computation gives $K(u)=Pu$; hence $Pu=0$ and
$u=S0=1$.

If a presentation uses a description named $W$ with defining equations
\[
 W(0)=1,\qquad W(Sx)=x,
\]
then it proves $K(u)=W(u)$ for every $u$ already from coverage.
In the zero branch both sides are $1$; in the branch $u=S(Pu)$ both
sides are $Pu$.  Thus $K$ and $W$ are provably equal on the whole domain,
even if they are distinct descriptions.  This equality does not appeal to
standard-model agreement or to uniqueness of an arbitrary recurrence.

\subsection{Four elementary uses of equation induction}

Choose descriptions for addition and truncated subtraction, with recursion
on the second argument:
\begin{align}
 x+0&=x, &x+Sy&=S(x+y),\label{eq:eic-addition}\\
 d(x,0)&=x, &d(x,Sy)&=P(d(x,y)).\label{eq:eic-subtraction}
\end{align}
Using these equations and \eqref{eq:eic-induction}, we now derive
\begin{align}
 0+y&=y,\label{eq:eic-left-zero}\\
 (Sx)+y&=S(x+y),\label{eq:eic-left-successor}\\
 d(Sx,Sy)&=d(x,y),\label{eq:eic-cancel-successors}\\
 d(Sa,a)&=1.\label{eq:eic-diagonal}
\end{align}
Each induction below has a single equation as its induction formula.

For \eqref{eq:eic-left-zero}, induct on $y$.  The base is $0+0=0$.
The step is $0+Sy=S(0+y)=Sy$.

For \eqref{eq:eic-left-successor}, hold $x$ fixed and induct on $y$.
At zero both sides are $Sx$.  At a successor the induction hypothesis
gives
\[
 (Sx)+Sy=S\bigl((Sx)+y\bigr)
       =S\bigl(S(x+y)\bigr)=S(x+Sy).
\]
Combining \eqref{eq:eic-left-successor} with the defining equation for
addition gives the useful rearrangement
\begin{equation}
 (Sy)+r=y+Sr.
 \label{eq:eic-shift-successor}
\end{equation}
Indeed, both sides are $S(y+r)$.

For \eqref{eq:eic-cancel-successors}, hold $x$ fixed and induct on $y$.
The base computes as
\[
 d(Sx,S0)=P(d(Sx,0))=P(Sx)=x=d(x,0).
\]
Applying $P$ to the induction equation gives the step:
\[
 d(Sx,S(Sy))=P(d(Sx,Sy))=P(d(x,y))=d(x,Sy).
\]

Finally, induct on $a$ for \eqref{eq:eic-diagonal}.  The base is
$d(S0,0)=S0=1$.  At a successor,
\[
 d(S(Sa),Sa)=d(Sa,a)=1,
\]
using \eqref{eq:eic-cancel-successors} and the induction hypothesis.

\subsection{An equation invariant that reflects actual equality}

The remaining arithmetic ingredient is the following equation, valid for
all $x,y$:
\begin{equation}
 C\bigl(\operatorname{sg}(d(x,y)),\ x,\ y+d(x,y)\bigr)=x.
 \label{eq:eic-residual-invariant}
\end{equation}
Its meaning is that a nonzero subtraction residual reconstructs the
original argument by addition.  The selector expresses this assertion as
one literal equation, so equation induction applies directly.

Hold $x$ fixed and induct on $y$.  At $y=0$, subtraction gives residual
$x$, and \eqref{eq:eic-left-zero} gives $0+x=x$.  The desired equation
is consequently $C(\operatorname{sg}(x),x,x)=x$, which is
\eqref{eq:eic-selector-idempotence}.

For the step, assume \eqref{eq:eic-residual-invariant} at $y$ and
abbreviate
\[
 r=d(x,y),\qquad r'=Pr=d(x,Sy).
\]
These are local term abbreviations, not additional function symbols.
Split coverage for $r'$.

If $r'=0$, then $\operatorname{sg}(r')=0$, and the successor instance
of \eqref{eq:eic-residual-invariant} reduces immediately to $x=x$.

In the other branch, $r'=S(Pr')$, so constructor separation gives
$r'\ne0$.  Split coverage for $r$.  Its zero branch would imply
$r'=P0=0$, a contradiction, from which the desired step equation follows.
In its successor branch,
\[
 r=S(Pr)=Sr',\qquad \operatorname{sg}(r)=1.
\]
The induction equation therefore computes to $y+r=x$.  Meanwhile,
\eqref{eq:eic-shift-successor} gives
\[
 (Sy)+r'=y+Sr'=y+r=x.
\]
Since $r'$ is a successor, $\operatorname{sg}(r')=1$, so this is
exactly the branch selected in the successor instance of
\eqref{eq:eic-residual-invariant}.  This completes the equation induction.
All case distinctions came from coverage; no induction on an implication
or a disjunction was used.

\paragraph{Reflection at residual $1$.}
We have, for arbitrary $a,b$,
\begin{equation}
 d(Sa,b)=1\quad\longleftrightarrow\quad a=b.
 \label{eq:eic-residual-reflection}
\end{equation}
For the forward direction, substitute $x=Sa$, $y=b$ into
\eqref{eq:eic-residual-invariant}.  If $d(Sa,b)=1$, the sign is $1$
and the selector returns its second branch, giving
\[
 b+1=Sa.
\]
The defining equation for addition gives $b+1=S(b+0)=Sb$.  Thus
$Sb=Sa$, and applying $P$ gives $b=a$, hence $a=b$.
Conversely, $a=b$ and \eqref{eq:eic-diagonal} give
$d(Sa,b)=d(Sa,a)=1$.

Define the equality characteristic by the fixed description
\begin{equation}
 e(a,b)=\operatorname{sg}\bigl(K(d(Sa,b))\bigr).
 \label{eq:eic-equality-characteristic}
\end{equation}
The sign, countdown, and residual reflection lemmas give
\begin{align}
 e(a,b)=0&\quad\longleftrightarrow\quad a=b,
       \label{eq:eic-equality-reflection}\\
 e(a,b)=0&\ \lor\ e(a,b)=1.
       \label{eq:eic-equality-range}
\end{align}
In detail, $e(a,b)=0$ is equivalent to $K(d(Sa,b))=0$ by
\eqref{eq:eic-sign-zero}, which is equivalent to $d(Sa,b)=1$ by
\eqref{eq:eic-countdown-zero}, and then to $a=b$ by
\eqref{eq:eic-residual-reflection}.  The range assertion is
\eqref{eq:eic-sign-range} applied to $K(d(Sa,b))$.

The reflection property concerns the logic's actual equality throughout
the domain.  It was not assumed as part of a numerical comparison symbol,
and its proof did not restrict $a$ or $b$ to standard numerals.
It also yields decidable equality.  In the $e(a,b)=0$ branch,
\eqref{eq:eic-equality-reflection} gives $a=b$.  In the $e(a,b)=1$
branch, an assumption $a=b$ would give $e(a,b)=0$ and therefore $0=1$,
contradicting separation.  Hence
\begin{equation}
 \forall a\,\forall b\,\bigl(a=b\lor(a=b\to\bot)\bigr)
 \label{eq:eic-equality-decidable}
\end{equation}
is a theorem, not an extra hypothesis.

\subsection{Compilation of all quantifier-free formulas}

For each fixed quantifier-free formula $A(\vec p)$, construct a term
$b_A(\vec p)$ by external recursion on its finite syntax.  Take equality,
$\bot$, $\land$, $\lor$, and $\to$ as the primitive formula constructors;
negation is implication into $\bot$.  Set
\begin{align}
 b_{s=t}&=e(s,t), & b_{\bot}&=1,\label{eq:eic-boolean-atoms}\\
 b_{A\land D}&=C(b_A,b_D,1),\label{eq:eic-boolean-and}\\
 b_{A\lor D}&=C(b_A,0,b_D),\label{eq:eic-boolean-or}\\
 b_{A\to D}&=C(b_A,b_D,0).\label{eq:eic-boolean-implies}
\end{align}
Every term here is a finite composition of fixed primitive-recursive
descriptions; if desired, projections and composition give it a single
description symbol.  No internal quantification over formulas or programs
is involved.

An external induction on the syntax of $A$ proves, within the arithmetic
theory,
\begin{equation}
 \bigl(b_A=0\lor b_A=1\bigr)
 \quad\text{and}\quad
 \bigl(b_A=0\longleftrightarrow A\bigr).
 \label{eq:eic-boolean-correctness}
\end{equation}
Here are the full connective checks.  In every case $b_A=1$ implies
$\neg A$: if $A$ held, the already established implication
$A\to b_A=0$ would contradict separation.

\paragraph{Atomic formulas and falsity.}
For an equation, both claims are
\eqref{eq:eic-equality-reflection}--\eqref{eq:eic-equality-range}.
For falsity the range assertion is $1=0\lor1=1$, with its right disjunct
proved by reflexivity.  Also $1=0\to\bot$ follows from separation and
symmetry, while $\bot\to1=0$ follows by explosion.  Thus $1=0$ is
equivalent to $\bot$, as required.

\paragraph{Conjunction.}
Split $b_A=0\lor b_A=1$.  In the zero branch, $A$ holds and
$b_{A\land D}=b_D$.  Its range follows from the range of $b_D$,
and its being zero is equivalent to $D$, hence to $A\land D$ under
the known $A$.  In the one branch, $\neg A$ holds and
$b_{A\land D}=1$.  Its range is immediate.  If it were zero,
separation would give $\bot$, hence $A\land D$; conversely an
assumption $A\land D$ contradicts $\neg A$, and explosion gives
$b_{A\land D}=0$.

\paragraph{Disjunction.}
In the branch $b_A=0$, the formula $A$ holds and
$b_{A\lor D}=0$.  The range is immediate, and both $b_{A\lor D}=0$
and $A\lor D$ hold, proving their equivalence.  In the branch
$b_A=1$, we have $\neg A$ and $b_{A\lor D}=b_D$.  Its range follows
from the range of $b_D$.  If its value is zero, then $D$ and hence
$A\lor D$.  Conversely, from $A\lor D$, the $A$ alternative
contradicts $\neg A$ and yields $b_D=0$ by explosion, while the
$D$ alternative yields $b_D=0$ directly.  This proves the required
equivalence in both alternatives.

\paragraph{Implication.}
In the branch $b_A=0$, the formula $A$ holds and
$b_{A\to D}=b_D$.  Its range follows from the range of $b_D$.
If its value is zero, then $D$, which gives $A\to D$.  Conversely,
$A\to D$ and the known $A$ give $D$, hence $b_D=0$.
In the branch $b_A=1$, we have $\neg A$ and
$b_{A\to D}=0$.  Its range is immediate.  To prove $A\to D$,
assume $A$; this contradicts $\neg A$, and explosion gives $D$.
Thus both $b_{A\to D}=0$ and $A\to D$ hold in this branch, and
their equivalence follows.  This argument uses the proved two-value
range of the finite characteristic.  It does not use excluded middle
for arbitrary formulas or change the surrounding intuitionistic logic.

This completes the external syntactic induction proving
\eqref{eq:eic-boolean-correctness} for every quantifier-free formula.

\subsection{Recovery of induction and the full theory}

Let $A(n,\vec p)$ be any fixed quantifier-free formula.  Suppose an
arbitrary retained context $\Gamma$ proves its base and its induction
step:
\[
 \Gamma\vdash A(0,\vec p),
 \qquad
 \Gamma,A(z,\vec p)\vdash A(Sz,\vec p),
\]
with the usual freshness condition on $z$.  By
\eqref{eq:eic-boolean-correctness}, the base gives
$\Gamma\vdash b_A(0,\vec p)=0$.  Under the additional assumption
$b_A(z,\vec p)=0$, the same equivalence gives $A(z,\vec p)$,
the given step gives $A(Sz,\vec p)$, and the equivalence gives
$b_A(Sz,\vec p)=0$.  Apply literal-equation induction to
$b_A(n,\vec p)=0$ and then apply its reflection implication.
The result is $\Gamma\vdash A(t,\vec p)$ for every target term $t$.
Thus the theory proves the full contextual quantifier-free induction
rule, including in arbitrary quantified contexts.

\paragraph{Completion theorem.}
In the common language of fixed primitive-recursive descriptions, the
theory
\[
 \mathsf{PRE}+\mathsf{Cov}+\mathsf{Sep}+\mathsf{EI}
\]
has exactly the same sequents as ordinary first-order intuitionistic
primitive recursive arithmetic with all quantifier-free induction
instances.  The forward simulation follows from the recovered
quantifier-free induction, successor injectivity, and successor separation;
all primitive-recursive defining equations and all logical rules are
shared.  For the reverse simulation, literal-equation induction is an
instance of quantifier-free induction, and successor separation gives
$\mathsf{Sep}$.  Quantifier-free induction proves $\mathsf{Cov}$:
at zero choose its left disjunct, and at a successor use
$Sx=S(P(Sx))$, which follows from $P(Sx)=x$, for its right disjunct.

Both interpretations therefore use the full domain, defined by $x=x$,
and preserve every function symbol, equality, logical connective, and
quantifier.  Both composites have the identity isomorphism, with graph
$x=y$.  The theories are consequently identity-domain bi-interpretable
and equiconsistent.  This completion theorem is unconditional once the
literal-equation induction rule has been obtained; it has no dependency
on a classical successor-uniqueness reduction.

\section{Pairing, primitive recursion, and internal numerical composites}
\label{sec:iterator-coding}
This appendix is adapted with attribution from the progress-closure
arithmetic development~\cite{progress}. Its prerequisites have now been
proved: scalar predecessor, arithmetic reflection, and full quantifier-free
induction. Pairing is a derived proof tool; it is not part of any checker.

\subsection{Pairing from a bounded square-root search}

Once QFI is available, a familiar square-based injection gives a
particularly direct compiler. Write $a\le b$ for $d(a,b)=0$.
The usual linear-order and semiring identities follow by fixed
quantifier-free inductions from the successor, addition, and monus
equations. Multiplication itself is the scalar term
\[
                        a\cdot b=I(b,0;u.u+a).
\]
Only primitive-recursive arithmetic is involved in the following
bounded root search; there is no minimization over an unbounded domain.

For a fixed parameter $n$, put
\[
 U_n(v)=C(d((Sv)^2,n),Sv,v),\qquad
 R_n(i)=I(i,0;v.U_n(v)),\qquad q(n)=R_n(n).
\]
The update increments $v$ exactly when $(v+1)^2\le n$, and
otherwise retains it. QFI in $i$, with $n$ fixed, proves
\[
 R_n(i)\le i,\qquad R_n(i)^2\le n,\qquad
 R_n(i)=i\ \vee\ n<(S(R_n(i)))^2.              \tag{Root invariant}
\]
The base has $R_n(0)=0$. If the update increments, its test
contradicts the strict-upper-bound alternative in the third invariant.
Thus the old state equals $i$, and the new state equals $Si$;
its square bound is the increment test. If the update does not increment, the
strict upper square bound holds and the state remains fixed.
An already established strict upper square bound forces every
subsequent update to remain fixed. These are exactly the cases
needed for the three displayed invariants.

At $i=n$, either the strict upper square bound already holds or
$R_n(n)=n$, in which case $n<(n+1)^2$ supplies it. Thus
\[
                        q(n)^2\le n<(q(n)+1)^2.           \tag{Root}
\]
Monotonicity of squares shows uniqueness of a number satisfying
Root. All formulas just used are quantifier free in the scalar
language, with fixed $n$ during the induction.

Now define
\[
 \pi(a,b)=(a+b)^2+a,\qquad
 L(n)=d(n,q(n)^2),\qquad H(n)=d(q(n),L(n)).
\]
For $k=a+b$, one has $k^2\le k^2+a<(k+1)^2$ since $a\le k$.
Root uniqueness gives $q(\pi(a,b))=a+b$, and therefore
\[
                    L(\pi(a,b))=a,\qquad H(\pi(a,b))=b.
\]
Surjectivity of $\pi$ is unnecessary. The original numerical domain
is never replaced by the range of this coding; pairing is used only
to store the finite state of a fixed primitive recursion. This
standard square pairing is a convenient reconstruction here, not
a claimed new arithmetic basis.

\subsection{Compiling each primitive recursion}

Suppose the fixed target description is
\[
 f(0,\vec p)=g(\vec p),\qquad
 f(Sn,\vec p)=h(n,f(n,\vec p),\vec p).
\]
Assume by external structural induction that the earlier $g,h$
already have scalar translations $g^*,h^*$. Define
\[
 W(u,\vec p)=\pi(SL(u),h^*(L(u),H(u),\vec p)),
\]
\[
 V(n,\vec p)=I(n,\pi(0,g^*(\vec p));u.W(u,\vec p)),
 \qquad f^*(n,\vec p)=H(V(n,\vec p)).
\]
The update has only its state variable bound, with every $p_i$
fixed. Equation induction proves $L(V(n,\vec p))=n$.
The two projection equations and scalar computation then prove
both defining equations for $f^*$ internally. Initial functions
and composition translate directly. This completes the translation
of every fixed finite PR description without adding a universal
evaluator.

\subsection{Why both numerical round trips are internal}

Two fixed scalar iterators with the same seed and provably equal
updates are internally equal at all counters: use literal-equation
induction comparing their values. Their base agrees, and their
successor equations plus congruence prove the step. The same
uniqueness argument works for ordinary primitive recursion.

For an ordinary PR description, induct externally on its finite
construction. The double translation satisfies its original
defining equations, so internal equation induction identifies it
with the original operation. For a native scalar description,
earlier translated updates agree by the external induction
hypothesis; scalar uniqueness identifies the double-translated
iterator. These proofs hold at arbitrary internal counters.

Thus the domain relation $x=y$ is provably a total bijection
preserving every fixed numerical operation in both composites.
The formula constructions and round trips in
Section~\ref{sec:choice-introduction} give the corresponding formula
composites. Both interpretations preserve falsity, taking native
$0=1$ to an ordinary contradiction and ordinary $\bot$ to native
$0=1$. Finite proof translations therefore establish equiconsistency.

\clearpage
\begin{thebibliography}{12}\raggedright
\bibitem{repo} V. Reshetnikov, ed., \emph{Concise presentations of intuitionistic
primitive recursive arithmetic}, ProveIt, commit
\texttt{11814771b7490dfbea0ed5fcf0f5b23586334764}.
\href{https://github.com/VladimirReshetnikov/ProveIt/tree/11814771b7490dfbea0ed5fcf0f5b23586334764/Logic/PrimitiveRecursiveArithmetic/ConciseIPRA}{Pinned repository snapshot}.
The README, filed packages and six incoming archives define the inspected
comparison frontier. Exact source paths and hashes are in
\texttt{SOURCE\_PROVENANCE.json}.
\bibitem{rawformal} ProveIt, \emph{Raw evaluation and uniform clauses},
\texttt{Lean/ConciseIPRA/NativeEval.lean},
\texttt{NativeClausal.lean}, and \texttt{NativeUniformClauses.lean},
in~\cite{repo}. In particular \texttt{eprf\_iff\_nprf},
\texttt{cprf\_iff\_nprf}, and \texttt{vprf\_iff\_cprf}.
\href{https://github.com/VladimirReshetnikov/ProveIt/blob/11814771b7490dfbea0ed5fcf0f5b23586334764/Logic/PrimitiveRecursiveArithmetic/ConciseIPRA/Lean/ConciseIPRA/NativeUniformClauses.lean}{Uniform-clause source}.
\bibitem{zeroformal} ProveIt, \emph{Seed evaluation and zero-normal clauses},
\texttt{Lean/ConciseIPRA/NativeZeroSeed.lean} and
\texttt{NativeZeroClausal.lean}, in~\cite{repo}.
\href{https://github.com/VladimirReshetnikov/ProveIt/blob/11814771b7490dfbea0ed5fcf0f5b23586334764/Logic/PrimitiveRecursiveArithmetic/ConciseIPRA/Lean/ConciseIPRA/NativeZeroSeed.lean}{SeedEval source};
\texttt{nvprf\_iff\_pra} and \texttt{zvprf\_iff\_nprf}.
\bibitem{choice} \emph{Complementary introductions with one local assumption},
\texttt{Research/certified\_zero/reference/inherited\_choice\_basis.tex},
in~\cite{repo}. Reproduced with attribution in this report.
\bibitem{bridge} \emph{Arithmetic strength and both ordinary interpretation
composites}, \texttt{Research/certified\_zero/reference/inherited\_ordinary\_bridge.tex},
in~\cite{repo}. Reproduced with attribution.
\bibitem{erasure} \emph{Internal scalar predecessor bootstrap},
\texttt{Research/certified\_zero/reference/scalar\_predecessor.tex},
in~\cite{repo}. Source of the factor map, comparison and saturated-iteration
construction of predecessor; reproduced with attribution.
\bibitem{forward} \emph{From literal-equation induction to full intuitionistic
PRA}, \texttt{Research/certified\_zero/reference/equality\_reflection.tex},
in~\cite{repo}. Internal reflection of actual equality and the compilation
of quantifier-free formulas; reproduced with attribution.
\bibitem{progress} \emph{Pairing, primitive recursion, and internal numerical
composites}, \texttt{Research/certified\_zero/reference/primitive\_recursion.tex},
in~\cite{repo}. Reproduced with attribution.
\end{thebibliography}
\end{document}
